\documentclass[11pt]{article}
\usepackage{amsmath,amssymb,amsthm}
\usepackage[margin=1.05in]{geometry}
\usepackage{microtype}
\usepackage{xurl}
\urlstyle{tt}
\emergencystretch=2.5em
\newtheorem{theorem}{Theorem}[section]
\newtheorem{proposition}[theorem]{Proposition}
\newtheorem{lemma}[theorem]{Lemma}
\newtheorem{corollary}[theorem]{Corollary}
\theoremstyle{definition}
\newtheorem{definition}[theorem]{Definition}
\newtheorem{hyp}{Input}
\renewcommand{\thehyp}{H\arabic{hyp}}
\theoremstyle{remark}
\newtheorem{remark}[theorem]{Remark}
\newcommand{\F}{\mathbb F}
\newcommand{\E}{\mathbb E}
\newcommand{\X}{\mathcal X}
\newcommand{\LL}{\mathcal L}
\newcommand{\DTIME}{\mathrm{DTIME}}
\newcommand{\NTIME}{\mathrm{NTIME}}
\newcommand{\coNTIME}{\mathrm{coNTIME}}
\newcommand{\parTIME}{\oplus\mathrm{TIME}}
\newcommand{\TRANS}{\mathrm{TRANS}}
\newcommand{\Ball}{\mathrm{Ball}}
\newcommand{\acc}{\#\mathrm{acc}}
\newcommand{\src}[1]{\url{#1}}
\title{Linear time with logarithmically many guesses is not in\\ deterministic time $n(\log n)^{\sigma}$\\[4pt]
\large (machine-checked in Lean 4, with no assumption)}
\author{Research driven, orchestrated and reviewed by Jack Ryan.\\
Claude (by Anthropic), GPT-6.2 Astra (by OpenAI) and Gemini 4 Neon (by Google)\\
collaborated in finding and writing the proofs.\\
Both theorems are machine-checked in Lean 4.\\
No professional mathematician has refereed the paper or the Lean proof.}
\date{25 November 2026; revised 26 November 2026\\
DOI: 10.5281/zenodo.24527297\\
Repository: \url{https://github.com/jackryan225/ntime-dtime-polylog}}
\begin{document}
\maketitle

\begin{abstract}
We prove the following statement about multitape Turing machines, and the proof is checked in Lean~4 with no assumption:
there is a constant $\sigma>0$ such that
\[
  \NTIME[n;\,2\log n]\ \not\subseteq\ \DTIME\bigl(n(\log n)^{\sigma}\bigr),
\]
where $\NTIME[n;g]$ is linear nondeterministic time with $g$ nondeterministic bits. In particular
$\NTIME(n)\not\subseteq\DTIME(n(\log n)^{\sigma})$. The constant $\sigma$ is not effective: the proof is a dichotomy on whether
characteristic $2$ ``has ports'' in the sense of the spectral geometry of tensors. If it has no ports, a char-$2$ form of the
OpenAI transposition lower bound gives $\Omega(n\sqrt{\log n})$ for deciding matrix transposition, whose complement uses
$\lfloor\log n\rfloor+1$ guess bits. If it has ports, the OpenAI iterated compressor, applied to a network whose base graph is the
coarse head-movement skeleton of a parity machine running in time $n^2$ and whose lanes are $\operatorname{polylog}$ many sub-blocks,
gives $\parTIME(n^2)\subseteq\oplus\cdot\forall\cdot\DTIME(n^2(\log n)^{-s}\operatorname{polyloglog}n)$ with a short
universal quantifier; together with the hypothesis to be refuted this contradicts the parity time hierarchy.
The proof of the second case contains a one-substep locality lemma for the compressor records (radius $7$ in the old base
graph), written out in full. We also prove the stronger statement
$\coNTIME[n;2\log n]\not\subseteq\NTIME(n(\log n)^{\sigma})\cap\parTIME(n(\log n)^{\sigma})$, so in particular
$\coNTIME[n;2\log n]\not\subseteq\mathrm{UTIME}(n(\log n)^{\sigma})$ (Section~\ref{sec:utime}). Under Weak ETH the first case alone
already gives the main statement, by a short argument from the OpenAI transposition lower bound (Section~\ref{sec:prior}); the
new part is the second case, which removes Weak ETH. The first version of this paper (25 November 2026) proved the theorem
relative to twelve stated inputs H1--H12 taken from three OpenAI preprints of 24 November 2026. The Lean proof assumes none
of them: it builds on OpenAI's Lean library and proves the rest (Section~\ref{sec:lean}). The inputs remain in this text as
cited results, and the two errors that the Lean work found in the first version are corrected.
\end{abstract}

\paragraph{Status.}
Theorem~\ref{thm:main} and Theorem~\ref{thm:Uplus} are proved in Lean~4 with no hypothesis. The Lean proofs depend only on the
three standard axioms of Lean and were checked by three kernels (Section~\ref{sec:lean}). This paper is the proof in words. It is
organised around twelve inputs H1--H12 (Section~\ref{sec:hyp}), which restate, with file and line, theorems and
constructions of the preprints \emph{Euclidean geometry in the asymptotic spectrum of tensors} (cited as EG~\cite{EG}), \emph{Circuit
lower bounds for Hamiltonian cycle and the permanent} (CL~\cite{CL}) and \emph{The logarithmic exponent of integer multiplication on tapes}
(T~\cite{T}), whose sources are in the public repository \src{github.com/openai/math}, folder \src{preprints/}. Every step of this
text that is not one of these inputs is proved here. In the first version the inputs were hypotheses of the theorem. They are
not hypotheses any more, because the Lean proof does not assume them. We restate them in our own words and copy only their formulas.
H4 is used in a characteristic-$2$ form that the preprints do not
state (we derived it from T's proof); H5 is stated in T for the rationals and for every field in OpenAI's Lean library; the price
comparison in H2 is stated in T for the rationals and for every field in CL and in OpenAI's Lean library; H3 is field-free; H1(d) is stated in EG's text for every field but formalized in OpenAI's Lean library
only in characteristic $0$. Section~\ref{sec:hyp} says exactly which change is made. The Lean work found two errors in the first version of this text
(one unchecked bit in check C12, and a missing length bound for the energy prefixes); both are corrected here
(Section~\ref{sec:lean}). Earlier versions of the argument were checked by blind re-derivations, a line-by-line check of
the port case, an adversarial reading and an independent reading of the whole paper; every repair they found is built in (the list
is in the file \src{CHECKS.md} that accompanies this paper). In particular the tag rule uses the full forbidden set (``Fixed
choices'' in Section~\ref{subsec:witness}), the degree is regularized as in EG (Section~\ref{subsec:network}(f)), head offsets ride on
the state wire, every padding field is checked to be zero, and the base-record cost is counted with the remaining lanes
(Lemma~\ref{lem:space}). That ports over $\F_2$ and the parameter solution in $\F_2$ follow from ports in characteristic $2$ was
checked separately (EG \texttt{lem:comparison}).

\tableofcontents

%======================================================================
\section{Introduction}\label{sec:intro}

Paul, Pippenger, Szemer\'edi and Trotter proved that linear nondeterministic time is strictly larger than linear
deterministic time for multitape Turing machines. The known separations have the form $\NTIME(n)\not\subseteq\DTIME(n\,f(n))$
with $f$ of iterated-logarithm type (see Section~\ref{sec:prior}). We prove a polylogarithmic separation.

\begin{theorem}[Main theorem]\label{thm:main}
There is $\sigma>0$ such that
\[
  \NTIME[n;\,2\log n]\not\subseteq\DTIME\bigl(n(\log n)^{\sigma}\bigr)
\]
for multitape Turing machines. Consequently $\NTIME(n)\not\subseteq\DTIME(n(\log n)^{\sigma})$.
\end{theorem}

Here $\NTIME[n;g]$ is the class of languages accepted in linear time by multitape machines that make at most $g(n)$ binary
guesses (Definition~\ref{def:classes}). Every such language is in $\DTIME(n^{3})$, so the separation is inside $\mathrm P$.
Theorem~\ref{thm:main} is unconditional: it is checked in Lean~4 with no hypothesis (Section~\ref{sec:lean}).

\paragraph{What is ours and what is OpenAI's.} The spectral geometry of tensors, the notion of ports, the transposition lower
bound for tapes and the iterated compressor are due to OpenAI (the three preprints EG, CL and T, and OpenAI's Lean library).
Under Weak ETH the main statement is a short corollary of OpenAI's transposition lower bound (Section~\ref{sec:prior}). Our
contribution is the following.
\begin{enumerate}
\item[(i)] The dichotomy on ports in characteristic~$2$, which removes Weak ETH. Case~1 is OpenAI's transposition argument with
Weak ETH replaced by ``characteristic $2$ has no ports''; it needs characteristic-$2$ forms of two steps of that argument (H4, H5).
\item[(ii)] The ports case (Case~2, Section~\ref{sec:case2}): the network whose base graph is the head-movement skeleton of a
parity machine, the Bene\v{s}-routed lanes, the locally checkable compressor records with the locality lemma
(Lemma~\ref{lem:locality}), the witness and its checks, and the wrapper with the parity time hierarchy (Section~\ref{sec:wrapper}).
\item[(iii)] The stronger form for nondeterministic, parity and unambiguous machines (Section~\ref{sec:utime}).
\item[(iv)] The Lean~4 proof of both theorems, built on OpenAI's Lean library (Section~\ref{sec:lean}).
\end{enumerate}

\paragraph{The dichotomy.} A field \emph{has ports} (H1) if bounded-degree regular graphs of unbounded order satisfy one
uniform nonlinear Poincar\'e inequality in all spectral wire metrics over that field. This depends only on the characteristic.
The proof splits on characteristic~$2$.
\begin{enumerate}
\item[\textbf{Case 1.}] \emph{Characteristic $2$ has no ports.} Then every deterministic multitape machine that decides matrix
transposition $\TRANS$ uses $\Omega(n\sqrt{\log n})$ steps on some input of each large length $n=2N^2$
(Proposition~\ref{prop:trans}). The complement of $\TRANS$ is in $\NTIME[n;\lfloor\log n\rfloor+1]$. Any $\sigma<1/2$ works
(Section~\ref{sec:case1}).
\item[\textbf{Case 2.}] \emph{Characteristic $2$ has ports.} Then there is $s>0$ with
$\parTIME(n^2)\subseteq\oplus[m]\cdot\forall[k]\cdot\DTIME(m\log\log n)$ for $m=n^2(\log n)^{-s+o(1)}$ and
$k=\log(n^2)+O(1)$ (Lemma~P, Section~\ref{sec:case2}). If $\NTIME[n;2\log n]\subseteq\DTIME(n(\log n)^{\sigma})$ with
$\sigma<s$, then the universal part collapses to deterministic time $o(n^2)$ after padding, so
$\parTIME(n^2)\subseteq\parTIME(o(n^2))$, contradicting the parity time hierarchy (Theorem~\ref{thm:hier}). Any
$\sigma<s$ works (Section~\ref{sec:wrapper}).
\end{enumerate}
Since we do not know which case holds, and $s$ depends on constants whose existence H8 asserts in Case~2 only, $\sigma$ is
not effective.

\paragraph{Why Case 2 does not lose a factor $\log t$.} A direct use of the compressor would route the data-dependent
head-movement skeleton through the lanes of a fixed lifted graph, at a cost of a factor $\Theta(\log t)$ in wire volume, which
the saving $2^{-T}\le(\log t)^{-s}$ of the compressor cannot pay. Here the skeleton enters only the \emph{base graph}, which the
compressor allows to be arbitrary (H8). Lanes are $N=(\log t)^{\Theta(1)}$ sub-blocks of a big tape block, which are also the
$N$ small time intervals of a big time interval; only inside a big interval are the lanes permuted, by a Bene\v{s} network of
depth $O(\log N)=O(\log\log t)$. The records that the compressor schedules (colours, tags, layouts, slot orders) are
functions of the summation variables and are checked locally with $o(t(\log t)^{-s})$ guessed bits, by the locality lemma
(Lemma~\ref{lem:locality}).

\paragraph{Organisation.} Section~\ref{sec:prior} compares with prior work. Section~\ref{sec:models} fixes machine
models, classes, the parity hierarchy and a batched-lookup lemma. Section~\ref{sec:hyp} states the inputs H1--H12 with their sources. Sections
\ref{sec:case1}, \ref{sec:case2} and \ref{sec:wrapper} prove the theorem. Section~\ref{sec:utime} proves the stronger form for
nondeterministic, parity and unambiguous machines. Section~\ref{sec:lean} describes the Lean proof: what is formalized where, and what it takes from OpenAI's library.
Section~\ref{sec:remarks} discusses the range of the exponent, the weakest
points and what is not claimed.

%======================================================================
\section{Prior work}\label{sec:prior}

\paragraph{Linear nondeterministic versus deterministic time.} Paul, Pippenger, Szemer\'edi and Trotter~\cite{PPST} proved that
some language is recognized in linear time by a nondeterministic machine with two work tapes and by no deterministic multitape
machine in linear time: $\DTIME(n)\ne\NTIME(n)$. Santhanam~\cite{San} proved $\DTIME(n\sqrt{\log^*n})\ne\NTIME(n\sqrt{\log^*n})$,
with the same time bound on both sides. The bound $\NTIME(n)\not\subseteq\DTIME(o(n(\log^*n)^{1/4}))$ is stated by Chen and
Wang~\cite[footnote 3]{CW} as the best they know; they attribute it to~\cite{PPST} and to an answer of R.~Williams on MathOverflow
(question 4953, 16 December 2009) that derives it from the proof of~\cite{PPST}. All these factors are of iterated-logarithm type. Theorem~\ref{thm:main} gives a polylogarithmic factor. For simulations in small space much more is known, for example $\mathrm{NTIBI}[n,n^{\varepsilon}]\not\subseteq
\mathrm{DTISP}[n^{c_\varepsilon},n^{o(1)}]$ (Williams~\cite[Theorem 1.6]{Wil}; in his notation $\mathrm{NTIBI}[t,b]$ is nondeterministic
time $t$ with $b$ nondeterministic bits and $\mathrm{DTISP}[t,s]$ is deterministic time $t$ and space $s$ simultaneously); such bounds restrict the space of the deterministic
machine, and Theorem~\ref{thm:main} does not.

\paragraph{The same classes under fine-grained hypotheses.} Lingas~\cite{Lin} studies exactly the classes $\mathrm{NDTIME}(f(n),n)$
of languages decided in time $O(n)$ by multitape machines with an $f(n)$-bit nondeterministic oracle, and proves separations from
fine-grained hypotheses (arXiv version 5): for example, if unweighted triangle detection needs time $N^{\omega-o(1)}$ then
$\mathrm{NDTIME}(\lceil\frac12\log n\rceil,n)\not\subseteq\DTIME(n^{\delta'/2})$ for every $\delta'<\omega$, and under the
low-dimensional orthogonal vectors conjecture $\mathrm{NDTIME}(\lceil\log n\rceil,n)\not\subseteq\DTIME(n^{2-\epsilon})$ for every
$\epsilon>0$; further forms use APSP, 3SUM, clique problems and ETH. Some of these premises have since
been reported false in preprints: ETH~\cite{OAI-ETH}, and the polynomial-range forms of 3SUM and APSP~\cite{AVW}; the triangle and
orthogonal-vectors premises and Lingas's large-weight forms are, as far as we know, open. Where his premises hold, his bounds are much
stronger than ours. Theorem~\ref{thm:main} assumes no fine-grained hypothesis, and no other hypothesis. Its proof in this text rests on results stated in the OpenAI
preprints (the inputs H1--H12), two of them (H4 and H5) in a characteristic-$2$ form that we derived from the proofs in T; the Lean
proof takes them from OpenAI's Lean library or proves what it needs (Section~\ref{sec:lean}).

\paragraph{What follows from the OpenAI preprints directly.} Under the variable-count Weak Exponential Time Hypothesis, T proves that
a deterministic multitape machine that writes the transpose of an $N\times N$ binary matrix on a $2^{-\rho N^2}$ fraction of the
inputs, for a fixed $0\le\rho<1-1/\sqrt2$, needs $\Omega(N^2\sqrt{\log N})$ steps (T Theorem \texttt{thm:transpose}, \src{T/lower.tex} l.236--250). Its network theorem
(T Theorem \texttt{thm:table}) applies equally to a machine that decides $\TRANS$ (Proposition~\ref{prop:trans}), and the
complement of $\TRANS$ is in $\NTIME[n;\lfloor\log n\rfloor+1]$ (Corollary~\ref{cor:case1}). So \emph{under Weak ETH}
$\NTIME[n;\lfloor\log n\rfloor+1]\not\subseteq\DTIME(n(\log n)^{\sigma})$ for every $\sigma<1/2$ is a short corollary of T. Our Case~1 is this
argument with Weak ETH replaced by ``characteristic $2$ has no ports''. (T uses Weak ETH to show that characteristic $0$ has no
ports and runs its argument over the rationals; that implication says nothing about characteristic $2$.) The new content of Theorem~\ref{thm:main} is Case~2, which
makes the combination free of Weak ETH. Likewise, OpenAI's Lean library proves its multiplication lower bound under Weak ETH also
for sound nondeterministic machines (\src{OAI/Computability/MultitapeWeakETH/Main.lean}); Lemma~\ref{lem:A} below is the analogue
for nondeterministic deciders of $\TRANS$ in characteristic~$2$. The dichotomy on ports is OpenAI's own template for circuits
(CL \src{intro.tex} l.131--139); we found it applied to uniform time classes nowhere.

\paragraph{Tightness of Case~1.} The reduction of Harvey and van der Hoeven~\cite[Corollary 6.2]{HvdH} decides transposition in
time $O(M(m)+m\log\log m)$ for a matrix of $m$ bits, here $m=N^2=n/2$, where $M(m)$ is the time of a multitape multiplier on
$m$-bit operands; T gives
$M(m)=m(\log m)^{1/2+o(1)}$ and uses this reduction itself (\src{T/lower.tex} l.344--356). So
$\TRANS\in\DTIME(n(\log n)^{1/2+o(1)})$, and Proposition~\ref{prop:trans} is tight up to a factor $(\log n)^{o(1)}$.
(Tal~\cite{Tal} gives $o(N\log N)$ sorting and a transposition bound from the September 2026 OpenAI multiplication preprint; we do
not need it.)

\paragraph{Anthropic's VP $\ne$ VNP paper.} Anthropic~\cite{Ant} (hexagonmath.org 2611.00954v1, 23 November 2026, with Lean) proves
arithmetic-circuit lower bounds $m^{c(\log m)^{r}}$ for $\mathrm{HC}_m$ over every field and for $\mathrm{PERM}_m$ when $p\ne2$
(Thm.~\texttt{s:allfields}). It uses a different dichotomy, on prices over the parties $\mathrm{Mat}_\mu(\F_2)$: a
\emph{contracting} hypothesis, which yields a Poincar\'e inequality for all prices on the unitriangular groups $U_h$ (whose
generating sets grow with $h$), versus a \emph{flat} hypothesis, under which Euclidean metrics are prices (Def.~\texttt{s:cases},
Thm.~\texttt{s:alt}). The paper does not show the two hypotheses exclusive. They are not ``ports'' versus ``no ports'': flatness
does not by itself exclude ports: Euclidean placements satisfy the expander Poincar\'e inequality, so the Euclidean prices that
flatness provides do not violate it.
In its contracting case it cuts the weighted runs of a tape automaton (a queue-tape model) into pieces (Prop.~\texttt{prop:pieces},
Thm.~\texttt{thm:cut}). This happens at a scale $v\ge2^{\sqrt{\log T}}$, with pieces of $\Theta(KT/k)$ bits, $k\le v^{64}$, that
are evaluated by an \emph{accelerator} derived from a hypothetical small circuit; a self-referential program gives $V=1+V$
(Thm.~\texttt{thm:noacc}). It contains no statement about multitape Turing time classes, no transposition bound, and no parity-time
speed-up. We found that our Theorem~\ref{thm:main} and its strengthening in Section~\ref{sec:utime} are not contained in it,
and we see no short derivation of them from it. The shared feature is the ``either case'' template followed by a diagonal over exact sums of weighted runs.
(Full-text check of the TeX source, 25 November 2026.)

\paragraph{Searches.} Three independent searches for earlier statements (25 November 2026, between 06:56 and 08:00 UTC; together
they covered arXiv, ECCC, GitHub, Zenodo, vibemathed.com, hexagonmath.org in full text, nilradical.ai, proofatlas.ai, the Palomar
registry, the Kourovka notebook, StackExchange, blogs, and the papers citing~\cite{PPST}, \cite{San} and the three OpenAI preprints)
found no earlier statement of Theorem~\ref{thm:main}, or of any $\NTIME(n)\not\subseteq\DTIME(n\,g(n))$ with $g\gg\log^*n$ that
does not assume a fine-grained hypothesis; the third search also covered Theorem~\ref{thm:Uplus} and found nothing. The closest
recent work we found is a subexponential deterministic simulation of multitape nondeterminism~\cite{SS} (an upper bound, consistent
with ours). This records searches; it is not a priority claim.

\paragraph{Tools.} We use the two-tape simulation of Book, Greibach and Wegbreit~\cite{BGW} in its parsimonious form, the time
hierarchy method~\cite{HS65,Cook73,SFM78,Zak83}, rearrangeable networks~\cite{Benes,Waksman}, and, for comparison only, the
block-respecting simulation of Hopcroft, Paul and Valiant~\cite{HPV}, which we avoid. For a prime $p$, $\mathrm{Mod}_p$ classes
are closed under complement~\cite[Lemma 6.6]{BG}; an oracle relative to which deterministic and nondeterministic linear time
coincide is given by Hemmerling~\cite[Proposition 15]{Hem} (see also~\cite[Corollary 1]{Hem13}).

%======================================================================
\section{Models, classes and two tools}\label{sec:models}

All logarithms are to base $2$. For a real $y$, $\lceil y\rceil$ and $\lfloor y\rfloor$ are the usual integer parts.

\subsection{Machines}
A \emph{$k$-tape machine} has $k\ge1$ tapes with cells $0,1,2,\dots$, one head per tape, a finite control and a finite tape
alphabet $\Gamma$ containing a blank and a left endmarker $\vdash$. At the start every tape holds $\vdash$ in cell $0$ and blanks
elsewhere, except tape $1$, which holds the input $x\in\{0,1\}^n$ in cells $1,\dots,n$; all heads are on cell $0$. A transition
reads the $k$ scanned symbols and the state, writes $k$ symbols, moves each head by $-1,0,+1$ and changes the state. We use the
syntactic convention that no transition writes $\vdash$ outside cell $0$, overwrites $\vdash$, or moves a head left while it
scans $\vdash$. (Every machine is converted to this form with a constant factor in time and the same number of accepting
computations.) Consequently no computation ever moves a head to a negative cell.

A \emph{deterministic} machine has at most one transition for each state and scanned tuple; a \emph{nondeterministic} machine
(NTM) has finitely many. A \emph{computation} on $x$ is a finite sequence of configurations, each obtained from the previous by
one transition, ending in a halting state; it is \emph{accepting} if this state is accepting. We write $\acc(M,x)$ for the
number of accepting computations. A machine \emph{runs in time} $f$ if on every input of length $n$ every computation halts
within $\max(f(n),1)$ steps. A function $f$ is \emph{time constructible} if some deterministic machine on input $1^n$ writes $f(n)$ in
binary in $O(f(n))$ steps; $n^2$ is time constructible.

\begin{definition}[Classes]\label{def:classes}
Let $f,g,m,k:\mathbb N\to\mathbb N$. In guess budgets, $\log n$ means $\log\max(n,2)$.
\begin{enumerate}
\item $\DTIME(f)$: languages decided by a deterministic multitape machine that halts within $C\max(f(n),1)$ steps on every input of length $n$, for a constant $C$ (this is time $O(f)$; the $\max$ with $1$ matters only where $f(n)=0$, e.g.\ $n\le1$ for $f(n)=n(\log n)^\sigma$, where the literal bound $0$ would allow only languages that are constant on these inputs).
\item $\NTIME[f;g]$: languages $A$ for which there are a deterministic multitape machine $V$ and $c>0$ such that $V$, given
  $x\in\{0,1\}^n$ on tape $1$ and $y\in\{0,1\}^{g(n)}$ on tape $2$, halts within $c\max(f(n),1)$ steps (for every $y$), and
  $x\in A\iff\exists y\in\{0,1\}^{g(n)}:\ V(x,y)=1$. $\coNTIME[f;g]$ is the same with $\forall$ (equivalently, the
  complements of the languages in $\NTIME[f;g]$).
\item $\NTIME(f)$: languages $A$ for which some NTM $M$ runs in time $O(f)$ and $x\in A\iff\acc(M,x)\ge1$. If $g=O(f)$ and a machine that reads its
  input $x$ can compute $g(|x|)$ within $O(f(|x|))$ steps (true for $f(n)=n$ and every budget used here), then $\NTIME[f;g]\subseteq\NTIME(f)$: the NTM
  writes $g(n)$ guessed bits on tape $2$, one binary branching per bit, and runs $V$.
\item $\parTIME(f)$: languages $A$ for which some NTM $M$ runs in time $O(f)$ and $x\in A\iff\acc(M,x)$ is odd.
\item $\mathrm{UTIME}(f)$: languages $A$ for which some NTM $M$ runs in time $O(f)$, has $\acc(M,x)\le1$ on every input, and
  $x\in A\iff\acc(M,x)=1$. Clearly $\DTIME(f)\subseteq\mathrm{UTIME}(f)\subseteq\NTIME(f)\cap\parTIME(f)$.
\item $\oplus[m]\cdot\forall[k]\cdot\DTIME(f)$: languages $A$ for which there is a deterministic multitape machine $V$ that halts
  within $O(f(n))$ steps on every triple $(x,u,\gamma)$ with $|x|=n$, $|u|=m(n)$, $|\gamma|=k(n)$ (written on three tapes), and
  \[
     x\in A\iff \#\{u\in\{0,1\}^{m(n)}:\ \forall\gamma\in\{0,1\}^{k(n)}\ V(x,u,\gamma)=1\}\ \text{is odd}.
  \]
\end{enumerate}
\end{definition}
$\DTIME(f)$ is closed under complement; hence $\NTIME[f;g]\subseteq\DTIME(h)$ implies $\coNTIME[f;g]\subseteq\DTIME(h)$.

\subsection{Counters and batched lookup}
\begin{lemma}[Countdown]\label{lem:countdown}
A multitape machine can count down a binary counter from $g$ to underflow in $O(g+\log g)$ steps.
\end{lemma}
\begin{proof}
Keep the counter on its own tape with the low bit under the head. A decrement flips the trailing zeros to ones and the next
one to zero, then returns. A bit at position $i$ is touched in at most $\lceil g/2^{i}\rceil+1$ of the $g$ decrements, so the
total head travel is $O(\sum_i (g2^{-i}+1))=O(g+\log g)$; underflow is the borrow running off the top bit.
\end{proof}

\begin{lemma}[Batched lookup]\label{lem:batch}
Let a string $w$ of length $\mu$ be on a tape, and let a work tape hold $r$ requests $(a_i,\ell_i)$ with $a_i+\ell_i\le\mu$,
written in canonical binary, one request per delimited record. A multitape machine writes the substrings $w[a_i,a_i+\ell_i)$, each tagged with its request number, in
$O\bigl(\mu+(r(\log(\mu+2)+\log(r+2))+\sum_i\ell_i)\log(r+2)\bigr)$ steps.
\end{lemma}
\begin{proof}
Tag each request with its number and sort the requests by $a_i$ (merge sort on two work tapes; a record has
$O(\log(\mu+2)+\log(r+2))$ bits, so $O(r(\log(\mu+2)+\log(r+2))\log(r+2))$ steps). Then scan $w$ from the left: after request $i$
has been copied the head on $w$ is at $a_i+\ell_i$, and it moves by $a_{i+1}-(a_i+\ell_i)$ to the next start, counting with
Lemma~\ref{lem:countdown}; this move is negative only if request $i+1$ starts inside request $i$, and then it is at most $\ell_i$
cells long. So the head on $w$ travels at most $\mu+2\sum_i\ell_i$ cells, and copying costs $O(\sum_i\ell_i)$ more. Sort the
answers, tagged with their numbers, back by request number ($O((\sum_i\ell_i+r\log(r+2))\log(r+2))$ steps).
\end{proof}

\subsection{The parity time hierarchy}

\begin{theorem}[Parity time hierarchy]\label{thm:hier}
Let $t$ be time constructible with $t(n)\ge n$. There is a language in $\parTIME(t)$ that is not in $\parTIME(t')$ for any $t'$
with $t'(n)=o(t(n))$.
\end{theorem}
\begin{proof}
\emph{Transcripts.} Let $M$ be a $k$-tape NTM with transition set $\Delta$. A \emph{transcript} is a word over $\Delta$. Fix a
machine $U$ with a fixed number of tapes that, given the code $\langle M\rangle$ and input $x$, (i) writes nondeterministically a
transcript $\delta_1\delta_2\cdots\delta_\ell$, one transition at a time, each as a fixed-length binary code (invalid codes
reject), stopping after the first transition that enters a halting state; (ii) checks that $\delta_1$ starts in the initial state
of $M$ and that the states in consecutive transitions agree, and for $i=1,\dots,k$ deterministically checks tape $i$: it replays
tape $i$ on a work tape, checking for every $j$ that $\delta_j$ reads the symbol currently under the replayed head; (iii) accepts
iff all checks pass and $\delta_\ell$ enters an accepting state. Each computation of $M$ on $x$ has exactly one transcript, and a
transcript passes all checks iff it is the transcript of a computation; the checks are deterministic. Hence the accepting
computations of $U$ on $(\langle M\rangle,x)$ are in bijection with those of $M$ on $x$, and if $M$ runs in time $t_M$, then every
computation of $U$ that writes the transcript of a computation of $M$ halts within $c_M\,t_M(n)$ steps, where $c_M$ depends only on
$\langle M\rangle$ (this is the two-tape simulation of Book, Greibach and Wegbreit, which is parsimonious). A computation of $U$
that writes an invalid transcript may run longer, but it never accepts.

\emph{Diagonalisation.} Let $D$ be the NTM that on input $z$ of length $n$ first writes $t(n)$ on a clock tape and then branches
into two computations. The first accepts at once. The second reads $z$ as $\langle M\rangle 0 1^j$ (rejecting if $z$ has no such
form), runs $U$ on $(\langle M\rangle,z)$ and decrements the clock once per step of $U$ (Lemma~\ref{lem:countdown}); it rejects on
underflow and otherwise accepts iff $U$ accepts. $D$ runs in time $O(t)$. Let $M$ run in time $t'=o(t)$. For all large $j$, on
$z=\langle M\rangle01^j$ every computation of $U$ on a valid transcript ends before the clock, and every other one rejects (by a
check of $U$ or by the clock), so $\acc(D,z)=1+\acc(M,z)$, which has the other parity.
Thus $L_\oplus(D)$ differs from $L_\oplus(M)$ at $z$. Every language of $\parTIME(t')$ is $L_\oplus(M)$ for such an $M$.
\end{proof}

%======================================================================
\section{The inputs H1--H12 from the OpenAI preprints}\label{sec:hyp}

Paths are relative to the folder \src{preprints/} of the public repository \src{github.com/openai/math}; Lean paths are relative to its
folder \src{lean/} (commit \texttt{3344da5}). Write
\src{EG/} for \src{Euclidean-geometry-in-the-asymptotic-spectrum-of-tensors-November-24-2026/build/},
\src{CL/} for \src{Circuit-lower-bounds-for-Hamiltonian-cycle-and-the-permanent-November-24-2026/build/} and
\src{T/} for \src{The-logarithmic-exponent-of-integer-multiplication-on-tapes-November-24-2026/build/}.
Each input is our restatement of the cited lines: formulas are copied, the prose is ours. Where the source says
``rational'' or ``Weak ETH'' and we use characteristic $2$, the change is stated in the input itself.
In the first version of this paper these twelve statements were hypotheses of Theorem~\ref{thm:main}. They are not hypotheses
any more: the Lean proof assumes none of them (Section~\ref{sec:lean}). They stay here as the interface between this text and
the preprints. Sections~\ref{sec:case1} to~\ref{sec:utime} are written relative to them, so the proof in words is exactly as
strong as H1--H12; the unconditional status of the theorems comes from the Lean proof.

\subsection{Spectral prices and ports}

\begin{hyp}[EG Definition \texttt{def:ports} and its discussion; \src{EG/port-graphs.tex} l.11--55]\label{H1}
Let $G$ be a symmetric stochastic matrix on a finite set $V$, and $x=(x_i)_{i\in V}$ points of a spectral space over a field
$\F$. Put
\[
 \mathcal E_G(x)=\frac1{|V|}\sum_{i,j\in V}G_{ij}d(x_i,x_j)^2,\qquad
 \mathcal V(x)=\frac1{|V|^2}\sum_{i,j\in V}d(x_i,x_j)^2 .
\]
(a) \emph{Definition.} $\F$ \emph{has ports} if there are a family of regular symmetric multigraphs of unbounded order and
bounded degree and a constant $\gamma$ such that $\mathcal V(x)\le\gamma\,\mathcal E_G(x)$ for every graph of the family (with
$G$ its normalized adjacency matrix, a diagonal arc keeping the walk in place) and every assignment $x$ into a spectral space
over $\F$ (l.20--29).
(b) Whether $\F$ has ports depends only on the characteristic of $\F$ (l.30--31).
(c) In (a) one may test either finite spectral points or complete $\mathrm{CAT}(0)$ spectral spaces (l.32--33).
(d) If $\F$ has ports, the family may be taken connected, loop-free, of one fixed power-of-two degree and of arbitrarily large
power-of-two orders (l.33--55).
\end{hyp}
We say \emph{characteristic $2$ has ports} if $\F_2$ has ports; by (b) this is the same for every field of characteristic $2$.

\begin{hyp}[The class $\X_2$; \src{CL/prices.tex} l.107--113 and l.129--149, and \src{T/graphs.tex} l.14--50]\label{H2}
Let $\X_2$ be the class of complete $\mathrm{CAT}(0)$ (Hadamard) spaces that carry compatible spectral prices for tensors over a
field of characteristic $2$ on all finite subsets, with distance equal to the price of a one-bit wire; it depends only on the
characteristic (CL l.129--141). A class of metric spaces \emph{admits ports} if regular multigraphs of a single positive degree
and unbounded orders satisfy one inequality $\E\,d(x_U,x_V)^2\le C\,\E\,d(x_I,x_J)^2$ ($U,V$ independent uniform vertices,
$(I,J)$ a uniform directed arc, diagonal arcs allowed) for all placements in all spaces of the class; for $\X_\chi$ this is
precisely port admission by the characteristic $\chi$ in the sense of H1 (CL l.142--149). Moreover (CL l.107--113, where
the comparison is stated without restriction on the field: applying the local coefficient tables of a network to the product of
its wire tensors supplies the output by restriction; T l.40--50 states it for rational tensors; OpenAI's Lean library states it
for every field as \texttt{wire\_cost\_le}, \src{OAI/Computability/MultitapeWeakETH/Geometry.lean} l.12): if the vertices of a
tensor network over $\F_2$ with edge widths $q_e$ are placed at
points $x_u$ of a member of $\X_2$, then any product of equality wires obtained from the network by local linear maps at the
vertices has price at most $\sum_{e=uv}q_e\,d(x_u,x_v)$, and $r$ copies cost $r$ times this.
\end{hyp}

\begin{remark}\label{rem:noports}
If characteristic $2$ has no ports, then $\X_2$ is a class of Hadamard spaces without ports: given any family of regular graphs
of one degree and any $C$, H1(a) fails for it with $\gamma=C$ at some assignment into a spectral space over $\F_2$; by H1(c)
this assignment can be taken in a complete $\mathrm{CAT}(0)$ spectral space, i.e.\ in a member of $\X_2$, and
$\mathcal E_G=\E\,d(x_I,x_J)^2$, $\mathcal V=\E\,d(x_U,x_V)^2$.
\end{remark}

\subsection{Inputs for Case 1}

\begin{hyp}[CL Theorem \texttt{thm:loggraph}; \src{CL/entropy.tex} l.54--66]\label{H3}
Let $P$ be the walk on a $d$-regular multigraph with $n\ge2$ vertices, $d\ge1$, and uniform stationary law $\pi$. If
$\mathcal V_\pi(x)\le\beta\,\mathcal E_P(x)$ for every placement $x$ in every member of a class of Hadamard spaces without ports,
then $\beta\ge c\,\log n/\log(d+1)$ for an absolute constant $c>0$. (Here $\mathcal E_P(x)=\sum_{i,j}\pi_iP_{ij}d(x_i,x_j)^2$
and $\mathcal V_\pi(x)=\sum_{i,j}\pi_i\pi_jd(x_i,x_j)^2$, l.14--19.)
\end{hyp}

\begin{hyp}[T Lemma \texttt{lem:geometry} in characteristic $2$; \src{T/lower.tex} l.42--116]\label{H4}
Assume characteristic $2$ has no ports. Let a multigraph have $P_0$ vertices, maximum degree $D$ and edge widths at most $Db$;
let $P$ be the least power of two $\ge P_0$. For all sufficiently large $P_0$ there are a placement $u\mapsto x_u$ in some
$X\in\X_2$ and a point $o\in X$ such that, with $S=\sum_{e=uv}q_e\,d(x_u,x_v)$ and constants depending only on $D$,
\begin{align}
 S &\le C_D\,bP/\sqrt{\log P}, \label{eq:short}\\
 \E_{a\in A,\,t\in B}\bigl(d(x_a,x_t)-1\bigr)^2 &\le C_D\,P/(N\log P), \label{eq:term}\\
 \textstyle\sum_v h_v\,d(x_v,o) &\le \tfrac1{\sqrt2}\textstyle\sum_vh_v+C_D\bigl(\tfrac{P}{\log P}\sum_vh_v^2\bigr)^{1/2}, \label{eq:central}
\end{align}
for any two lists $A,B$ of $N$ vertices, distinct within each list (independent uniform draws), and all weights $h_v\ge0$.\\
\emph{Change from the source:} T assumes Weak ETH and the class $\X$ for rational tensors. Its proof uses Weak ETH only to
conclude, via T Theorem \texttt{thm:noports}, that $\X$ is a class of Hadamard spaces without ports (l.70--73), and then applies
H3; everything else (expander padding, scalar Poincar\'e inequality, barycenter in a $\mathrm{CAT}(0)$ space, Cauchy--Schwarz)
is field-free. We replace that step by Remark~\ref{rem:noports}. This transfer was audited by two independent readings of T's proof.
\end{hyp}

\begin{hyp}[T Lemma \texttt{lem:entropy} over $\F_2$; \src{T/lower.tex} l.130--202]\label{H5}
Let a tensor network over $\F_2$ have pairs of vertices $(a_j,t_j)$, $1\le j\le L$, and let $W=(W_1,\dots,W_L)$ be a random
binary vector. Suppose that for each $w$ in the support of $W$, pins $C_v(w)$ at the vertices make the contracted tensor over
$\F_2$ equal to
\[
 \bigotimes_{j=1}^L\bigl((e_{w_j})_{a_j}\otimes(e_{w_j})_{t_j}\bigr).
\]
Then for every placement in $X\in\X_2$, every $o\in X$ and every ordering of the coordinates,
\begin{equation}\label{eq:entropy}
 \sum_{j=1}^{L}H(W_j\mid W_1,\dots,W_{j-1})\,d(x_{a_j},x_{t_j})\ \le\ S+\sum_vH(C_v(W))\,d(o,x_v).
\end{equation}
\emph{Change from the source:} ``over the rationals'' and ``nonzero scalar multiple'' become ``over $\F_2$'' and ``equal to''
(the only nonzero scalar of $\F_2$ is $1$). In the proof (l.147--202) the field enters only through the price comparison of H2
and the division by $c_*=\prod_wc(w)^{rp(w)}$, which is $1$; the tensor placed at $o$ is the $0/1$ indicator of the joint pin
tuples of the selected arrays, and distinct colour tuples give distinct arrays and hence distinct joint pin tuples, so no
coefficient is a sum of two ones. This transfer was audited by two independent readings of T's proof. Moreover OpenAI's Lean library states this lemma for every field: \texttt{paired\_entropy\_at\_placement} and
\texttt{paired\_entropy} (\src{OAI/Computability/MultitapeWeakETH/EntropyPrice.lean} l.164 and l.125, under \texttt{[Field K]}); with
$K=\F_2$ they give H5. This was compared at the level of statements; we did not build these files.
\end{hyp}

\subsection{Inputs for Case 2: the compressor of EG}

Throughout this subsection $\F$ is a field; we shall use $\F=\F_2$.

\begin{hyp}[EG Theorem \texttt{thm:conversion}; \src{EG/compressor.tex} l.16--31 and l.36--63]\label{H6}
Fix a finite set of at least two parties, a multiset $D$ of demand edges and a connected spanning multigraph $S$ of supply edges,
each edge with distinct endpoints. For integer widths $b=(b_e)_{e\in D}\ge0$ let $W_D(b)$ be the product of identity tensors of
dimensions $2^{b_e}$ on the edges; define $W_S(c)$ likewise. Give the edges positive rational masses $p_e$ ($e\in D$), $q_e$
($e\in S$), and put $\|b\|_p^2=\sum_{e\in D}b_e^2/p_e$, $\|c\|_q^2=\sum_{e\in S}c_e^2/q_e$. For a spectral point over $\F$ on the
parties write $d_e$ for its wire distance on the endpoints of $e$. Let $H\ge0$ and assume that for every such spectral point
\[
 \textstyle\sum_{e\in D}p_ed_e^2\le H\sum_{e\in S}q_ed_e^2 .
\]
Then for every $H'>H$ and all $b$ there is an ordinary sum identity
\[
 W_D(b)=\sum_\alpha\Bigl(\bigotimes_vL_{v,\alpha}\Bigr)W_S(c_\alpha),\qquad \|c_\alpha\|_q^2\le H'\|b\|_p^2+O(1),
\]
with the constant independent of $b$. The identities use finitely many scalar parameters $\kappa$ subject to fixed integer
polynomial equations with a solution in $\F$, and hold over every commutative unital ring with any assignment satisfying them.
With $M=\max(\{0\}\cup\{b_e\})$, each index $\alpha$ is a fixed number of integers of magnitude $O(M+1)$; admission and supply
widths are computable in $\mathrm{poly}(\log(2+M))$ bit operations; given $\alpha$ and its incident bit strings, an entry of
$L_{v,\alpha}$ is zero or a monomial in $\kappa$ of degree $O(M+1)$, with a description computable in time polynomial in $M+2$.
\end{hyp}

\begin{hyp}[Wire networks and local identities; \src{EG/compressor.tex} l.267--312]\label{H7}
A \emph{network of the stated form} is: a \emph{base}, a finite indexed directed multigraph with distinct endpoints on every arc
and with $d$ incoming and $d$ outgoing arcs at every vertex, $d$ a power of two, the slots at each vertex numbered separately in
each direction; a displacement $\tau_e$ for each base arc $e:x\to y$ in a binary vector space of \emph{lanes}; positions $(x,g)$
($x$ a base vertex, $g$ a lane); actual arcs $(x,g)\to(y,g+\tau_e)$, each of an integer width $b\ge0$ that may vary with the
lane; an actual arc of width $b$ equates two strings of $b$ bits; with $P$ positions there are $Pd$ actual arcs, each of mass
$p=1/(Pd)$, and $\LL=\sum_{\text{actual}}b_e^2/p$, so $\sum b_e\le\sqrt\LL$. Arbitrary vertex factors sit at the positions.
A \emph{substep} gives an identity
\begin{equation}\label{eq:kernel}
 W((a_u)_u)=\sum_\iota\sum_z\prod_vK_{v,\iota}\bigl((a_u)_{\pi_\iota(u)=v},\,z|_v\bigr),
\end{equation}
where $W$ is the product of the wire equalities, $a_u$ the incident half-wire indices at $u$, $\iota$ a list of admitted motif
terms that determines the new topology, widths, the placement $\pi_\iota$ of old positions and the kernels, and $z$ one string
per new actual arc, used at both of its ends. A kernel consists of local ties, sums over local indices and entries of the maps of
H6. After multiplication by arbitrary old vertex factors and summation over the $a_u$, every whole old factor belongs to the new
vertex $\pi_\iota(u)$, with its uncontracted indices. The identities hold in every commutative unital ring satisfying the fixed
equations.
\end{hyp}

\begin{hyp}[EG Theorem \texttt{thm:compressor}; \src{EG/compressor.tex} l.314--340, with l.524--527 and l.548--551 of its proof]\label{H8}
Suppose $\F$ has ports. There are a power-of-two degree $d_0$, an integer $b_0\ge1$, a constant $C_2\ge2$ and finitely many motif
templates, all fixed, such that every network of the stated form (H7) with $P_0$ positions of degree $d_0$, maximum width $W_0$
and at least $Tb_0$ lane coordinates, $T\ge1$, has an exact expansion by $T$ blocks of local identities \eqref{eq:kernel},
ending at degree $d_0$. One finite system of parameter equations, satisfiable in $\F$, serves every network and every $T$. For
every sequence of admitted terms,
\begin{equation}\label{eq:compress}
 \textstyle\sum_{e\ \mathrm{final}}b_e\le\sqrt{\LL_T}\le2^{-T}\bigl(\sqrt{\LL_0}+O(P_0)\bigr).
\end{equation}
All intermediate wire and kernel fields have length at most $C_2^T(W_0+1)$. At each final position the recursive expansion of
its vertex tensor uses at most $C_2^T$ occurrences of vertices, map entries and fields (each initial factor and each map entry
counted as one unevaluated operation, one occurrence shared by all values of its hidden indices); for fixed arguments a nonzero
map entry is a monomial of parameter degree at most $C_2^T(W_0+1)$. Moreover (proof, l.524--527): each block consists of a fixed
number $s_0$ of shrinking substeps (H10) followed by one refresh substep (H11); each shrinking substep divides the number of
positions by its own fixed $K$, the refresh keeps it, and the product $\kappa_0>2$ of the $K$'s of a block gives
$P_i=P_0/\kappa_0^{i}$ positions at the start of block $i$; and (proof, l.548--551) once final arc bits and term choices are fixed,
each final vertex sums only its own hidden indices, because each substep sends each entire old vertex to one new position.
\emph{Scope.} The expansion is the one built by H9--H11, and all of the above holds for every proper colouring in H9(i), every
choice of tag offsets that satisfies H9(v), every division of payloads as evenly as possible and either tie rule of the layout:
EG's proof uses only these properties of the choices.
\end{hyp}

\begin{hyp}[EG Lemma \texttt{lem:compressor-slots} and the construction in its proof; \src{EG/compressor-appendix.tex} l.9--105]\label{H9}
Let the current base degree be $d$ (a power of two), $K=2^{2d+1}d$, $Q=3K/4$. A shrinking substep uses the following packing.
\begin{enumerate}
\item[(i)] Colour the base greedily in base order with $2d+1$ colours, so that the endpoints of every base arc get distinct
  colours (l.29, l.99--100).
\item[(ii)] Split the lane as $(U,h,j,g)$: $U\in\F_2^{2d+1}$ has one bit per colour, $h$ indexes the $d$ outgoing slots, $j$
  is a tag in a binary vector space of order $J$, $g$ is the remaining lane (l.29--32). A position $(x,U,h,j,g)$ is
  \emph{thick} if $U_{c(x)}=1$ and \emph{thin} otherwise (l.32--33).
\item[(iii)] For given $x,j,g$ all thick positions are packed into one new position with base $(x,j)$ and lane $g$; there are
  $P/K$ new positions. The $K/4$ thick--thick wires of each base arc, for fixed $j,g$ at its source, are concatenated in a common
  order at both ends into a \emph{bundle} from a packed position $A$ to a packed position $B$. Every other wire has a thin end,
  taken as $C$, and the other end as $D$; a thin--thin wire is taken in its old direction $C\to D$ (l.39--45).
\item[(iv)] Each base vertex $x$ has a tag offset $\phi_x$. A thin position $(x,U,h,j,g)$ is assigned to the bundle with source
  $(x,j+\phi_x,g)$ on the outgoing base arc of slot $h$ at $x$; every wire on which it plays $C$ is assigned to that bundle, as a
  $CD$ demand, and each bundle receives exactly $Q$ of them (l.47--54).
\item[(v)] With $L_x=\{0\}\cup\{\tau^{\rm tag}_f: f\text{ leaves }x\}$, the offsets satisfy, at every incidence of a base arc $e$
  between $x$ (the base of $C$) and $y$ (the base of $D$, $e$ in either direction),
  $\phi_x\notin\tau^{\rm tag}_e+L_x$ and $\phi_x+\phi_y\notin\tau_e^{\rm tag}+L_x+L_y$; they are chosen greedily in base order,
  which works for every power of two $J>2d((d+1)+(d+1)^2)$ (l.56--72). Then $A,B$ and $D^*=\pi(D)$ have distinct new bases in each
  motif, in either layout.
\item[(vi)] Each of the roles $A,B,D^*$ occurs exactly $dQ$ times at every new position (l.74--85).
\item[(vii)] Packing, assignment and both choices of image depend on the remaining lane only through the arc translations. The
  parallel slots of a surviving occurrence are labelled by the original $CD$ arc identifying the motif, the surviving type and the
  subdivision index; a common deficit is filled by zero-width arcs to the next tag cyclically at the same old base and remaining
  lane; incoming and outgoing slots are numbered separately by sorting their identifiers, with all remaining lanes expressed
  relative to the vertex being numbered; widths and term-dependent map entries never enter the identifiers (l.87--97).
\end{enumerate}
\end{hyp}

\begin{hyp}[The shrinking substep; \src{EG/compressor.tex} l.344--445]\label{H10}
With the packing of H9: the payload of each bundle is divided into $Q$ integral pieces as evenly as possible, and each piece is
paired with one of the $Q$ demands $CD$ assigned to the bundle; the piece is the demand $AB$ of a six-party motif on
$A,B,C,D,M,N$ (l.361--366). Each motif is expanded by H6 with fixed masses; one term index per motif is chosen (l.374--381). The
supply groups are $G_1$ (the triangle on $A,C,M$ and the edge $DN$) and $G_2$ (the triangle on $B,C,N$ and the edge $DM$). A single
layout is chosen for the whole list of terms, retaining the group of smaller total energy with deterministic tie breaking
(l.388--392). In layout~1 a thick position goes to its packed site, a thin position to the $A$ of its assigned bundle, the $D$ map
of every motif to $D^*=\pi(D)$, the fresh $M$ to $A$ and the fresh $N$ to $D^*$, so $G_1$ becomes local and the surviving edges
$BC,BN,CN,DM$ join $B$--$A$, $B$--$D^*$, $A$--$D^*$, $A$--$D^*$; layout~2 interchanges $A,B$ and $M,N$ (l.392--404). Each
surviving type $e$ gets $m_e=\lfloor pq_e/p'\rfloor$ parallel arcs per occurrence ($p,p'$ the old and new arc masses, $q_e$ the
relative supply mass of $e$, the output degree taken large enough that $m_e>0$; l.410--412), oriented by its two roles in a fixed order, its payload split as evenly
as possible among them, and the common deficit is filled with zero-width arcs as in H9(vii); the slots commute with translations of
the remaining lane (l.405--422). The kernel at a new position concatenates and slices the thick--thick fields identically at both
ends, puts the maps of $C,D$ at the images of these old endpoints and all other maps by the layout, and sums each absorbed wire
inside its new position; a new vertex receives boundedly many old positions, maps and fields, each field of length
$O_{d,d'}(W+1)$; concatenations and slices are bounded lists of field identifiers, offsets and lengths (l.431--445).
\end{hyp}

\begin{hyp}[The refresh substep; \src{EG/compressor.tex} l.447--509]\label{H11}
The last shrinking substep of a block has output degree $D$, an order of the port family of H1(d), with port graph of degree $d_H$.
The refresh consumes $\log_2D$ lane coordinates as a port index $j$, included in the new base label; for each old base vertex
$x$ and remaining lane $g$ the $D$ ports form a block. An old base arc $e:x\to y$ maps the ports of a source block bijectively to
those of one destination block, by translation of port and remaining lane; these $D$ wires are concatenated in source-port order
into one external wire from the port numbered by the outgoing slot of $e$ at $x$ to the port numbered by its incoming slot at
$y$; every port has exactly one outgoing and one incoming external arc. Inside each block every old half-wire is tied to its
segment of the external wire at its designated port; the ties are supplied by the port-graph arcs of the block through one H6
term per block. Deficits are padded by zero-width arcs from each port to the next in cyclic order; all slots commute with
translations of the remaining lane; one old position stays at each new position.
\end{hyp}

\begin{hyp}[EG Lemma \texttt{lem:descriptors}; \src{EG/compressor.tex} l.561--666]\label{H12}
Under the hypotheses of H8 (with $n\ge P_0$, $W_0\le n^{O(1)}$, $T=O(\log n)$) there is a nondeterministic RAM procedure whose
successful paths are in bijection with the sequences of admissible term indices and which constructs at every level: for each
actual incident slot its width and both endpoint slots (base slots give the lane displacements); for each new position the list
of its immediate old positions, its local maps with their term data, and the slices, ties and summed fields of its kernel
\eqref{eq:kernel}. In its proof (l.598--660): term indices are guessed in a fixed order with fixed-length binary encodings,
out-of-range or inadmissible codes are rejected and every admitted index occurs once; the two candidate surviving energies are
accumulated and the layout is chosen by the prescribed tie rule, the common mass cancelling and the scores being sums of squared
integer widths with fixed rational coefficients; bounded lists are ordered by identifiers whose lane coordinates are written
relative to the bundle's source lane or the numbered position (by xor), so that the base records computed at lane $0$ are the base
records for every lane; a refresh lists, for each old base arc and source block, its $D$ actual arcs in source-port order.
\end{hyp}

\begin{remark}[Which inputs are used where]
Case~1 uses H1(a)--(c), H2, H3 (through H4), H4 and H5. Case~2 uses H1(a),(b),(d) and H6--H12. Lemma~\ref{lem:locality} uses
only the constructions in H9--H12, not their energy estimates. The energy estimate is used only through \eqref{eq:compress}.
\end{remark}

%======================================================================
\section{Case 1: characteristic 2 has no ports}\label{sec:case1}

Throughout this section assume that characteristic $2$ has no ports.

\subsection{Partial transposition over $\F_2$}
A \emph{tensor network over $\F_2$} is a multigraph whose edges carry bit strings of given widths, whose vertices carry tables
($\F_2$-valued functions of the incident strings and of the external indices at the vertex); the contraction fixes the external
indices, multiplies the tables and sums over all internal strings in $\F_2$. Edges between the same two vertices form a
\emph{bundle}.

\begin{proposition}[T Theorem \texttt{thm:table} over $\F_2$]\label{prop:table}
Fix $0\le\rho<1-1/\sqrt2$ and $D\ge1$. Let a tensor network over $\F_2$ have maximum degree $D$ in edge bundles, each bundle of
width at most $DN$, and $N$ distinct row terminals $A_1,\dots,A_N$ and $N$ distinct column terminals $B_1,\dots,B_N$, each
carrying one external index in $\{0,1\}^N$, and no other external index. Suppose that for a set $\mathcal G$ of at least
$2^{(1-\rho)N^2}$ matrices $X\in\{0,1\}^{N\times N}$, fixing the external index at $A_i$ to row $i$ of $X$ ($1\le i\le N$) makes
the contraction, as a tensor in the column indices, equal to $\bigotimes_k e_{\mathrm{col}_k(X)}$ at $B_k$. Then the network has
$P_0\ge c_{D,\rho}N\sqrt{\log N}$ vertices for all $N\ge N_{D,\rho}$.
\end{proposition}
\begin{proof}
This is the proof of T (l.208--230), with H4 and H5 in place of their sources. Write $n=N^2$ and suppose $P_0\le N\sqrt{\log N}$;
also $P_0\ge2N$. Modify each $A_i$: its table now takes a \emph{pin} $c\in\{0,1\}^N$, feeds $c$ into the network as the old external
index, and exposes $c$ again on a fresh external index of $N$ one-bit slots at $A_i$. Let $W$ be uniform on $\mathcal G$, viewed as
$L=n$ bits $W_{ik}$; pins are $C_{A_i}(W)=$ row $i$ of $W$, and no other vertex has a pin. For $w\in\mathcal G$ the pinned network
contracts to $\bigotimes_{i,k}\bigl((e_{w_{ik}})_{A_i}\otimes(e_{w_{ik}})_{B_k}\bigr)$ (slot $k$ at $A_i$, slot $i$ at $B_k$),
which is the hypothesis of H5 with pairs $(a_{(i,k)},t_{(i,k)})=(A_i,B_k)$. Take the placement and $o$ of H4 with $b=N$ and the
lists $A=(A_i)$, $B=(B_k)$. Since $P\le2P_0=O(N\sqrt{\log N})$, \eqref{eq:term} and Cauchy--Schwarz give
$\sum_{i,k}|d(x_{A_i},x_{B_k})-1|\le n\,(C_DP/(N\log P))^{1/2}=o(n)$; as conditional entropies lie in $[0,1]$, the left side of
\eqref{eq:entropy} is at least $H(W)-o(n)\ge(1-\rho)n-o(n)$ by the chain rule. On the right, \eqref{eq:short} gives
$S=O_D(NP/\sqrt{\log P})$, and \eqref{eq:central} with $h_{A_i}=H(\text{row }i)\le N$, $h_v=0$ otherwise
($\sum h_v\le n$, $\sum h_v^2\le N^3$) gives at most $n/\sqrt2+C_D(N^4/\sqrt{\log N})^{1/2}=n/\sqrt2+o(n)$. Hence
$(1-\rho-1/\sqrt2)\,n\le O_D(NP_0/\sqrt{\log P_0})+o(n)$, and since $1-\rho-1/\sqrt2>0$ this forces $P_0=\Omega_{D,\rho}(N\sqrt{\log N})$.
\end{proof}

\subsection{Transposition on tapes}
Let $\TRANS$ be the set of strings $w\in\{0,1\}^{2N^2}$, $N\ge1$, such that $w=XY$ with $X,Y\in\{0,1\}^{N\times N}$ in row-major
order and $Y=X^{T}$.

\begin{proposition}\label{prop:trans}
Assume characteristic $2$ has no ports. For every deterministic multitape machine $M$ deciding $\TRANS$ there are $c_M>0$ and
$N_M$ such that for every $N\ge N_M$ some input $(X,X^T)$ makes $M$ run at least $c_M\,N^2\sqrt{\log N}$ steps.
\end{proposition}
\begin{proof}
This follows T Theorem \texttt{thm:transpose} (\src{T/lower.tex} l.236--342) with two changes: the field is $\F_2$, and
$M$ decides a language rather than writing an output. Fix $N$ and let $T$ be the maximum running time of $M$ on the inputs
$(X,X^T)$; assume $T\le N^2\sqrt{\log N}$ (otherwise we are done). The input occupies cells $1,\dots,2N^2$ of tape $1$. Divide every
tape into \emph{blocks} of $N$ cells, block $j$ being the cells $[1+jN,1+(j+1)N)$ (cell $0$ is a block of its own); so row $i$ of
$X$ is block $i$ and row $k$ of $Y$ is block $N+k$ of tape $1$. Divide time into \emph{epochs} of $N$ steps, freezing a halted
configuration, up to the first multiple of $N$ that is $\ge T$.

\emph{Records.} The record of a run is the list of the blocks scanned by each head at each epoch boundary. A head moves at most
one block per epoch, so there are at most $2^{O_M(1+T/N)}$ records. Let $R$ be the record of the largest number of runs on inputs
$(X,X^T)$ and $\mathcal G_R$ the set of those $X$: $|\mathcal G_R|\ge2^{N^2-O_M(1+T/N)}\ge2^{(1-\rho)N^2}$ with $\rho=1/4$, for
$N\ge N_M$, because $T/N\le N\sqrt{\log N}=o(N^2)$.

\emph{Network.} One vertex per epoch. During an epoch a head reaches only its starting block and the two neighbours. For each such
block carry its contents after the epoch to the next epoch whose window contains it; pass the state and the in-block offsets of
all heads to the next epoch. Tables: the $0/1$ indicator that the epoch's incoming data, run for $N$ steps by $M$, give the outgoing
data and the block changes prescribed by $R$. Feed each block its initial contents at its first use: blocks $i$ and $N+k$ of tape~1
receive them from terminals $A_i$ and $B_k$; all other initial contents ($\vdash$, blanks) are fixed in the tables. At the last
use, sum over the final contents. The last epoch's table also requires an accepting halted state. Block wires have width
$O_M(N)$, state wires $O_M(\log N)$, the maximum degree in bundles is $O_M(1)$, and there are $O_M(N+1+T/N)$ vertices.

\emph{Exactness over $\F_2$.} Fix $X$ at the $A_i$ and $Y$ at the $B_k$. A labelling of the internal edges with all tables
equal to $1$ is a sequence of $N$-step runs whose boundary data agree; concatenated in time order they form a run of $M$ on
$(X,Y)$ that follows $R$ and accepts by time $T$. Conversely the run of $M$ on $(X,Y)$, if it follows $R$ and accepts by $T$,
gives such a labelling. Since $M$ is deterministic, each epoch table is a function of its inputs, and by induction over epochs in
time order there is at most one such labelling. Hence the contraction equals $[M$ accepts $(X,Y)$ by $T$ with record $R]\in\{0,1\}$:
no cancellation modulo~$2$ can occur. For $X\in\mathcal G_R$ it is $1$ at $Y=X^T$ and $0$ at every other $Y$ ($M$ rejects them).
Since row $k$ of $X^T$ is column $k$ of $X$, the hypothesis of Proposition~\ref{prop:table} holds with $D=O_M(1)$. Therefore
$O_M(N+1+T/N)\ge c\,N\sqrt{\log N}$, so $T\ge c'_MN^2\sqrt{\log N}$ for large $N$.
\end{proof}

\begin{corollary}\label{cor:case1}
If characteristic $2$ has no ports, then $\NTIME[n;\lfloor\log n\rfloor+1]\not\subseteq\DTIME(n(\log n)^{\sigma})$ for every $\sigma<1/2$.
\end{corollary}
\begin{proof}
The complement $\overline{\TRANS}$ is in $\NTIME[n;\lfloor\log n\rfloor+1]$: on input $w$ of length $n$, count $n$ in binary
(Lemma~\ref{lem:countdown}), compute $N=\sqrt{n/2}$ if $n\ge2$ and it is an integer (otherwise accept; $\operatorname{poly}(\log n)$ steps of
arithmetic), read the guess $y$ as a pair $(i,k)$ of $\lceil\log N\rceil$-bit numbers (reject if $i\ge N$ or $k\ge N$), compute
the positions $iN+k$ and $N^2+kN+i$, fetch the two bits by one scan with countdowns (Lemma~\ref{lem:batch}), and accept iff they
differ. This takes $O(n)$ steps and $2\lceil\log N\rceil\le\lfloor\log n\rfloor+1$ guess bits, and $w\notin\TRANS$ iff some guess accepts.
If $\NTIME[n;\lfloor\log n\rfloor+1]\subseteq\DTIME(n(\log n)^\sigma)$ then $\TRANS\in\DTIME(n(\log n)^\sigma)$, and a machine running in time
$O(n(\log n)^\sigma)=o(n\sqrt{\log n})$ contradicts Proposition~\ref{prop:trans} at $n=2N^2$.
\end{proof}

%======================================================================
\section{Case 2: characteristic 2 has ports}\label{sec:case2}

Throughout this section assume that characteristic $2$ has ports, i.e.\ $\F_2$ has ports. Apply H8 with $\F=\F_2$ and fix, once
and for all, its constants $d_0,b_0,C_2$, the block structure ($s_0$, the degrees chosen in each block, $\kappa_0$), the port family
of H1(d), the finitely many templates and a solution $\kappa\in\F_2$ of the parameter equations. These are finite data; the machine
built below has them hardwired. Put $c_2=\lceil\log C_2\rceil$ and
\[
  s=2^{-\lceil\log(64(1+b_0+c_2))\rceil},\qquad\text{so}\qquad 0<s\le\frac{1}{64(1+b_0+c_2)}.
\]

\begin{lemma}[Lemma P]\label{lem:P}
Let $A\in\parTIME(n^2)$. For $n\ge n_M$ (a constant fixed below) put $t=n^2$, $\lambda=\lceil\log t\rceil$,
$p=\lfloor\log\lambda\rfloor$, $T=\lfloor sp\rfloor$,
\[
 m=(p+1)^6\,\lceil t/2^{T}\rceil\qquad\text{and}\qquad k=\lceil\log t\rceil ,
\]
and for $n<n_M$ put $m=k=1$ and $\lambda=2$. Then $A\in\oplus[m]\cdot\forall[k]\cdot\DTIME(m\log\lambda)$. Moreover
$t\lambda^{-s}(\log\lambda)^6\le m\le2^9\,t\lambda^{-s}(\log\lambda)^6$ for large $n$.
\end{lemma}
(The bounds on $m$: $2^{-T}\ge2^{-sp}\ge\lambda^{-s}$ and $p+1>\log\lambda$; and $2^{-T}\le4\lambda^{-s}$ by \eqref{eq:p3} below,
$p+1\le2\log\lambda$. We use this $m$ rather than $\lceil t\lambda^{-s}(\log\lambda)^6\rceil$ because it is an integer function of $n$
computable by integer arithmetic.)

The rest of this section proves Lemma~\ref{lem:P}. Fix an NTM $M$ with $r$ tapes and an integer $a_M$ such that $M$ runs in time
$a_Mt$ and $x\in A\iff\acc(M,x)$ is odd. Constants written $O_M(\cdot)$ or $c_M$ may depend on $M$ and on the fixed data above. It
suffices to treat $n\ge n_M$ for a constant $n_M$: for the finitely many shorter inputs the verifier accepts iff $u=0^m$ and
$x\in A$ (a finite table), which gives the right parity.

\subsection{Parameters}\label{subsec:params}
Put $p=\lfloor\log\lambda\rfloor$ and
\[
 B=2^{\lfloor p/2\rfloor},\qquad h=\lfloor 19p/2\rfloor,\qquad N=2^h,\qquad b=BN,\qquad T=\lfloor sp\rfloor,\qquad
 I=\lceil a_M2^{\lambda}/b\rceil .
\]
$B$ is the \emph{small block} length and the \emph{small interval} length, $N$ the number of lanes, $b$ the \emph{big block} and
\emph{big interval} length, $T$ the number of compressor blocks and $I$ the number of big intervals. For $n\ge n_M$:
\begin{align}
 &T\ge1,\qquad \lambda^{1/2}/4\le B\le\lambda^{1/2},\qquad Ib\ge a_Mt,\qquad I\,b\le 4a_Mt, \label{eq:p1}\\
 &h-b_0T\ \ge\ \tfrac{19}{2}p-1-\tfrac{p}{64}\ \ge\ 9p,\qquad\text{so}\quad 2^{h-b_0T}\ge(\lambda/2)^9, \label{eq:p2}\\
 &C_2^{T}\le2^{c_2sp}\le\lambda^{c_2s}\le\lambda^{1/64},\qquad 2^{-T}\le2^{1-sp}\le4\lambda^{-s}. \label{eq:p3}
\end{align}
(For \eqref{eq:p2}: $b_0T\le b_0sp\le p/64$. For \eqref{eq:p3}: $c_2s\le1/64$ and $p\ge\log\lambda-1$.)

\subsection{Normal form of $M$ and skeletons}\label{subsec:normal}
Give every step of $M$ a \emph{transition code} in $\{0,1\}^{e_M}$, with exactly one valid code for each transition available in
the current state and scanned tuple. In a halting configuration exactly one code is valid, and it leaves the configuration
unchanged. A \emph{run} is a sequence of $\tau=Ib$ valid codes together with the configurations they produce. Each computation of
$M$ on $x$ (length $\le a_Mt\le\tau$) extends to exactly one run, and the run ends in an accepting state iff the computation
accepts. So $\acc(M,x)$ equals the number of accepting runs.

For the network we index tape cells by $\mathbb Z$ and let negative cells hold blanks initially. By the endmarker convention no
run moves a head to a negative cell; this is how skeletons that ``leave the tape'' are handled: the network below is
defined for \emph{every} skeleton, and for a skeleton that a run cannot follow it contracts to $0$.

\emph{Small block} $j\in\mathbb Z$ of a tape is the cells $[jB,(j+1)B)$; \emph{big block} $Q$ is the small blocks
$[QN,(Q+1)N)$. \emph{Small interval} $i\in[0,IN)$ is the steps $[iB,(i+1)B)$; \emph{big interval} $\iota\in[0,I)$ is the small
intervals $[\iota N,(\iota+1)N)$. In $B$ steps a head moves at most $B$ cells, so its small-block index changes by
$\sigma_\rho(i)\in\{-1,0,1\}$ over small interval $i$ on tape $\rho$. The \emph{skeleton} of a run is
$\sigma=(\sigma_\rho(i))\in\{-1,0,1\}^{rIN}$. Put $\beta_\rho(i)=\sum_{i'<i}\sigma_\rho(i')$, the small block of head $\rho$ at
the start of small interval $i$. During small interval $i$ head $\rho$ stays in the \emph{window}
$\beta_\rho(i)+\{-1,0,1\}$.

\begin{lemma}[Ownership]\label{lem:own}
Let $q_\rho(\iota)=\lfloor\beta_\rho(\iota N)/N\rfloor$. Every small block in a window of a small interval of big interval
$\iota$ on tape $\rho$ lies in big block $q_\rho(\iota)-1$, $q_\rho(\iota)$ or $q_\rho(\iota)+1$.
\end{lemma}
\begin{proof}
Write $\beta_\rho(\iota N)=qN+\varrho$ with $0\le\varrho<N$. For $0\le g<N$, $|\beta_\rho(\iota N+g)-\beta_\rho(\iota N)|\le g$,
so the window of small interval $\iota N+g$ lies in $[qN+\varrho-N,\,qN+\varrho+N]\subseteq[(q-1)N,(q+2)N-1]$.
\end{proof}
We say $\iota$ \emph{owns} big blocks $q_\rho(\iota)+\delta$, $\delta\in\{-1,0,1\}$, on tape $\rho$. For a tape $\rho$ and a big
block $Q$, its \emph{owners} are the $\iota$ that own it, in increasing order.

\subsection{The network $N_{x,\sigma}$}\label{subsec:network}
Fix $x$ and any $\sigma\in\{-1,0,1\}^{rIN}$. All base arcs below have displacement $0$ unless stated. Lanes are $\F_2^h$; we
identify a lane $g$ with an integer in $[0,N)$.

\emph{(a) Interval tables.} One base vertex $F(\iota)$ per big interval; its position at lane $g$ is the small interval
$i=\iota N+g$. Its incident strings are: the \emph{state} in and out (state of $M$ and the offsets of all heads inside their
small blocks, $O_M(\log B)$ bits), and for each tape $\rho$ and $\delta\in\{-1,0,1\}$ the \emph{content} in and out of small block
$\beta_\rho(i)+\delta$ ($B\lceil\log|\Gamma|\rceil$ bits). Its factor $\Phi_i$ is the number modulo $2$ of code strings in
$(\{0,1\}^{e_M})^B$ that are valid at every step and lead from the incoming configuration restricted to the windows to the
outgoing one, with the block changes $\sigma_\rho(i)$.

\emph{(b) Ports and owner links.} For each $\iota,\rho,\delta$ there are base vertices $\mathrm{In}(\iota,\rho,\delta)$ and
$\mathrm{Out}(\iota,\rho,\delta)$; their lane $g$ carries small block $(q_\rho(\iota)+\delta)N+g$. For consecutive owners
$\iota_a<\iota_{a+1}$ of a big block $Q=q_\rho(\iota_a)+\delta=q_\rho(\iota_{a+1})+\delta'$ there is a base arc
$\mathrm{Out}(\iota_a,\rho,\delta)\to\mathrm{In}(\iota_{a+1},\rho,\delta')$ (lane $g$ to lane $g$). At the first owner the
factor of $\mathrm{In}$ at lane $g$ forces its outgoing string to be the initial contents of small block $QN+g$ (from $x$,
$\vdash$, blanks); otherwise it ties outgoing to incoming. At the last owner $\mathrm{Out}$ has no outgoing link and factor $1$
(its incoming string is summed).

\emph{(c) The access permutation.} Fix $\iota,\rho$ and write $q=q_\rho(\iota)$. The \emph{sources} are $\mathrm{In}(\delta)$ at
lane $g$ and the content output of window $\delta$ of small interval $\iota N+g$; the \emph{destinations} are
$\mathrm{Out}(\delta)$ at lane $g$ and the content input of window $\delta$ of small interval $\iota N+g$
($\delta\in\{-1,0,1\}$, $g\in[0,N)$). For each small block $Z\in[(q-1)N,(q+2)N)$ write $Z=(q+\delta_Z)N+g_Z$ and list in
increasing $g$ the pairs $(g_1,\delta_1),\dots,(g_a,\delta_a)$ with $\beta_\rho(\iota N+g_c)+\delta_c=Z$. Send
$\mathrm{In}(\delta_Z)$ at $g_Z$ to the input $(g_1,\delta_1)$, the output $(g_c,\delta_c)$ to the input $(g_{c+1},\delta_{c+1})$,
and the output $(g_a,\delta_a)$ to $\mathrm{Out}(\delta_Z)$ at $g_Z$; if $a=0$ send $\mathrm{In}(\delta_Z)$ at $g_Z$ to
$\mathrm{Out}(\delta_Z)$ at $g_Z$. By Lemma~\ref{lem:own} every window pair occurs in exactly one list, so this is a bijection
between the $6N$ sources and the $6N$ destinations. Add two dummy types to make $8$ types; dummies go to dummies and carry zeros.
An item has an address (type $\pi\in\F_2^3$, lane $g\in\F_2^h$).

\emph{(d) Routing.} Number the $8$ source types and, separately, the $8$ destination types by $\F_2^3$, so that the access bijection is a
permutation of $\F_2^{h+3}$. Route it by the Bene\v{s} network in butterfly form: position columns $\kappa=0,\dots,2(h+3)-1$, and
between columns $\kappa$ and $\kappa+1$ a stage of switches pairing the addresses $a,a+e$ for one address bit $e=e(\kappa)$
(the bits in the order $h+3,\dots,1,\dots,h+3$). The switch settings (straight: $a\mapsto a$; cross: $a\mapsto a+e$) are computed
by the deterministic looping algorithm (each cycle of the bipartite constraint graph is two-coloured starting with colour $0$ at its
least-labelled edge, then recursion on the two halves), in time $\mathrm{poly}(N)$. Position column $\kappa$ consists of the base
vertices $R(\iota,\rho,\kappa,\pi)$, $\pi\in\F_2^3$, with lanes $\F_2^h$. If $e(\kappa)$ is a lane bit $e_j$, there are two base arcs
$R(\kappa,\pi)\to R(\kappa+1,\pi)$ with displacements $0$ and $e_j$; if it is a type bit $e'_j$, there are base arcs
$R(\kappa,\pi)\to R(\kappa+1,\pi)$ and $R(\kappa,\pi)\to R(\kappa+1,\pi+e'_j)$, both with displacement $0$. So each position has
exactly the two out-arcs and the two in-arcs of its switches. The source of type $\pi$ feeds $R(0,\pi)$ and $R(2h+5,\pi)$ feeds the
destination of type $\pi$, lane to lane. The factor of a position copies the string on its selected in-arc (selected by the
setting of the stage before it; the source arc in column $0$) to its selected out-arc (by the setting of the stage after it; the
destination arc in the last column), puts $0$ on its other out-arc, and requires its unselected in-arc to carry $0$. These factors
depend on $\sigma$ through the settings.

\emph{(e) The state chain.} Per big interval, sources are an external input $E^{\rm in}(\iota)$ (lane $0$; other lanes dummy) and
the state outputs of the $N$ small intervals; destinations are the state inputs and an external output $E^{\rm out}(\iota)$
(lane $0$). Map $E^{\rm in}$ at $0$ to the input of $g=0$, the output of $g$ to the input of $g+1$, the output of $N-1$ to
$E^{\rm out}$ at $0$, and dummies to dummies. Route it by a Bene\v{s} network on $2^{h+1}$ addresses as in (d). Base arcs
$E^{\rm out}(\iota)\to E^{\rm in}(\iota+1)$. $E^{\rm in}(0)$ forces the initial state (state of $M$, all offsets $0$) at lane $0$;
$E^{\rm out}(I-1)$ requires an accepting state at lane $0$. Unused lanes are forced to $0$.

\emph{(f) Degree $d_0$.} The base built so far has bounded valence $V_M$. Replace each base vertex $v$ of valence $a$ by a central
site $v_0$ and a path $v_1,\dots,v_a$; the $c$-th incident arc of $v$ now ends at $v_c$ (same displacement); the path edges
$v_0$--$v_1$--$\cdots$--$v_a$ carry the tuple of all $a$ incident strings (width $O_M(B)$), each $v_c$ ($c\ge1$) equates its arc's
string with the $c$-th component and passes the tuple on, and $v_0$ applies the original factor to the tuple, lane by lane. For the
port vertices $\mathrm{In}$ and $\mathrm{Out}$ the path always has a site for the owner link of (b), also when the link is absent
(first or last owner), and $E^{\rm in}(0)$ and $E^{\rm out}(I-1)$ likewise have a site for the state link of (e); such a site then
carries, in place of the link, a width-$0$ arc to the twin introduced below, which counts toward its valence. So $a$, and the path,
depend neither on $\sigma$ nor on $\iota$. Give every undirected edge both orientations, the string on one and width $0$ on the other. Now every site has in- and
out-degree equal to its valence $\le3$. Add a twin $v'$ of every site with the same arcs among twins (width $0$, factor $1$) and add
$d_0-\mathrm{val}(v)$ arcs $v\to v'$ and as many $v'\to v$, of width $0$ ($d_0\ge4$ by H8 and EG l.451). Every base vertex now has
$d_0$ in- and $d_0$ out-arcs, no arc is a loop, and all new arcs have displacement $0$. Number the slots by a fixed rule.
This is the padding of EG \src{exact.tex} l.109--115.

The result $N_{x,\sigma}$ is a network of the form of H7: base $B_0$, degree $d_0$, lanes $\F_2^h$ ($h\ge Tb_0$ coordinates by
\eqref{eq:p2}). Its base vertex set is $[0,I)\times V_{\rm loc}$ for a fixed finite-times-$O(h)$ local set $V_{\rm loc}$ that does
not depend on $\sigma$ (by the convention for $\mathrm{In}$, $\mathrm{Out}$, $E^{\rm in}(0)$ and $E^{\rm out}(I-1)$ in (f)); only the endpoints of the owner-link slots
(another owner, or the twin) and the factors depend on $\sigma$ and $x$. Hence
\begin{equation}\label{eq:size}
 P_0=I\,|V_{\rm loc}|\,N=O_M(Ih N)=O_M\!\Bigl(\frac{t\log\lambda}{B}\Bigr),\qquad W_0=O_M(B),\qquad
 \sqrt{\LL_0}\le d_0P_0W_0=O_M(t\log\lambda),
\end{equation}
using $IN B=Ib\le4a_Mt$ and $\LL_0=P_0d_0\sum b_e^2\le(P_0d_0W_0)^2$.

\begin{lemma}[Path count]\label{lem:pathcount}
For every $\sigma\in\{-1,0,1\}^{rIN}$ the contraction $Z_\sigma(x)\in\F_2$ of $N_{x,\sigma}$ equals the number modulo $2$ of
accepting runs of $M$ on $x$ with skeleton $\sigma$. Hence $\sum_\sigma Z_\sigma(x)=\acc(M,x)\bmod2$.
\end{lemma}
\begin{proof}
Expanding each $\Phi_i$ as a sum over code strings, $Z_\sigma(x)$ is the number modulo $2$ of pairs (labelling of all strings,
code strings for all small intervals) at which every factor is $1$. Given such a pair, all factors in (b)--(f) are ties or forced
zeros, so every content input of small interval $i$ equals, by (c), the content output of the previous access to the same small
block inside the same big interval, or the string entering through the in-port, which by (b) is the last output in the previous
owner (or passes through owners that do not access the block), or the initial contents; and every state input equals the
previous state output (e). By Lemma~\ref{lem:own} no small block is accessed in a big interval that does not own its big block.
By induction over small intervals in time order, the incoming data of interval $i$ are the restriction to the windows of the
configuration reached after the code strings of intervals $<i$, and the code string of $i$ continues the run by $B$ steps with
block changes $\sigma(i)$. So the pair defines a run with skeleton $\sigma$, accepting by (e). Conversely an accepting run with
skeleton $\sigma$ determines its code strings and every string of the network (contents and states at the corresponding times;
routed and regularization strings are then forced), and this pair has all factors $1$. The two maps are inverse.
A run has exactly one skeleton, so summing over $\sigma$ gives $\acc(M,x)$ modulo $2$.
\end{proof}

\subsection{Compression}\label{subsec:compression}
Apply H8 to $N_{x,\sigma}$ with $T$ blocks, over $\F_2$ with the fixed $\kappa$. Composing the identities \eqref{eq:kernel} of
all substeps and multiplying by the vertex factors (H7, H8) gives, for every $\sigma$,
\begin{equation}\label{eq:expansion}
 Z_\sigma(x)=\sum_{\iota\in\mathrm{Adm}(\sigma)}\ \sum_{z\in\{0,1\}^{\ell(\sigma,\iota)}}\ \prod_{v}W_v(z|_v;\iota),
\end{equation}
where $\mathrm{Adm}(\sigma)$ is the set of sequences of admitted terms (one per motif of each shrinking substep and one per block
of each refresh), $\ell(\sigma,\iota)$ is the total final width, $v$ runs over the final positions and $W_v\in\F_2$ is the
vertex tensor of $v$: the sum over its own hidden indices of the product of the initial factors and map entries in its
recursive expansion (H8, last sentence). By \eqref{eq:compress}, \eqref{eq:p3} and \eqref{eq:size},
\begin{equation}\label{eq:zlen}
 \ell(\sigma,\iota)\le2^{-T}\bigl(\sqrt{\LL_0}+O(P_0)\bigr)\le c_M\,2^{-T}t\log\lambda\le 4c_M\,t\,\lambda^{-s}\log\lambda .
\end{equation}
By H8 and \eqref{eq:p3}, at one final position there are at most $C_2^T$ occurrences, each field has at most
$C_2^T(W_0+1)$ bits, so the hidden field bits number at most
\begin{equation}\label{eq:hidden}
 C_2^{2T}(W_0+1)\le c_M\lambda^{1/2+1/32}=c_M\lambda^{17/32}.
\end{equation}

\subsection{Levels, panels and addresses}\label{subsec:panels}
Fix $(\sigma,\iota)$ with $\iota\in\mathrm{Adm}(\sigma)$. The expansion passes through \emph{levels} $\ell=0,1,\dots,\Lambda$,
$\Lambda=T(s_0+1)$; substep $\ell$ (shrinking or refresh) turns the level-$\ell$ network into the level-$(\ell+1)$ network. Write
$B_\ell$ for the base at level $\ell$, $d_\ell$ for its degree (a fixed constant for each $\ell\bmod(s_0+1)$), $\mathrm{dist}_\ell$
for the distance in the underlying undirected graph of $B_\ell$ and $\Ball_\ell(x,\varrho)$ for the set of bases at distance at most
$\varrho$ from $x$. The number of lanes at every level is at least $2^{h-b_0T}\ge(\lambda/2)^9$ by \eqref{eq:p2}, and at most $N$.

\begin{definition}[Panel]\label{def:panel}
The \emph{panel} of a base $y\in B_\ell$ consists of:
\begin{enumerate}
\item[(P1)] the \emph{base record}: for each of the $d_\ell$ outgoing and $d_\ell$ incoming slots of $y$, the address of the
  neighbouring base, the slot number of the arc there, and the displacement;
\item[(P2)] the widths of the outgoing actual arcs at all positions $(y,g)$;
\item[(P3)] if substep $\ell$ is shrinking: the colour $c(y)$, the tag offset $\phi_y$, the term index of every motif whose
  \emph{anchor} (the base of its thin end $C$) is $y$, and the energy prefixes $E_1(y),E_2(y)$ defined in
  Section~\ref{subsec:witness}; if it is a refresh: the term index of each block $(y,g)$;
\item[(P4)] if $\ell=\Lambda$: the prefix $\Omega(y)$ of the total outgoing final width of all bases with smaller address.
\end{enumerate}
Every field has a fixed length: addresses and prefixes $\lambda+O(\log\lambda)$ bits, displacements $h+O(1)$ bits, widths and the
integers of a term index $O(\log\lambda)$ bits (since all widths are at most $C_2^T(W_0+1)=\lambda^{O(1)}$ and H6 bounds the
integers of an index by $O(\text{width}+1)$), colours and offsets $O(1)$ bits. The energy prefixes fit as well: a level has at most
$P_0=O_M(t)$ positions, a position anchors at most $2d_\ell$ motifs, and every supply width is at most
$C_2^T(W_0+1)=\lambda^{O(1)}$, so $0\le E_k(y)\le2^{\lambda+O(\log\lambda)}$, and a field of $\lambda+O(\log\lambda)$ bits holds every
energy prefix for $n\ge n_M$ (the threshold of Lemma~\ref{lem:space}). If a coefficient of an energy can be negative, the stored
number is the cleared-denominator form with non-negative coefficients. A position anchors at most $2d_\ell$ motifs, so the
panel has a fixed number of term slots per position; unused slots are zero. Thus a panel has a fixed size $\Pi_\ell$, laid out as
a per-base part ((P1), colour, offset, $E_1,E_2$, $\Omega$: $O(d_\ell\lambda)$ bits) followed by one per-position part for each lane
((P2) and the term slots: $O(d_\ell\log\lambda)$ bits), so
$\Pi_\ell=O(d_\ell\lambda)+O(d_\ell\log\lambda)\cdot(\#\text{lanes at level }\ell)$, which is $\lambda^{O(1)}$. These are the two
kinds of fields that Lemma~\ref{lem:space} counts.
\end{definition}

\begin{remark}[Addresses]
Level-$0$ bases have the product addresses $(\iota,v)\in[0,I)\times V_{\rm loc}$. A base created at a shrinking substep from the
old base $x$ and tag $j$ is $(x,j)$ (H9(iii)); one created at a refresh from $x$ and port $j$ is $(x,j)$ (H11). We give it the
product address $\mathrm{addr}(x)\cdot J_\ell+j$ (resp.\ $\mathrm{addr}(x)\cdot D+j$), where a tag $j$ is read as its binary
number as in (R2) (and a port as its number). Thus $B_{\ell+1}$ is the full product
$B_\ell\times[0,J_\ell)$ or $B_\ell\times[0,D)$, every address is in use, the panel of $y$ sits at the fixed offset
$\mathrm{off}_\ell+\mathrm{addr}(y)\,\Pi_\ell$ of the witness, and the parent of a base is its address divided by $J_\ell$ or $D$.
Addresses are therefore never stored for their own base; the neighbour addresses in (P1) are stored, and they are what lets the
verifier walk through a ball. This answers a point of the line-by-line check of the port case: the neighbour's address is recomputed from its
parent and tag (or port), and in addition stored in (P1) at $O(\lambda)$ bits per slot, which is counted in
Lemma~\ref{lem:space}.
\end{remark}

\subsection{The one-substep locality lemma}\label{subsec:locality}

\begin{lemma}[One-substep locality]\label{lem:locality}
Fix $(\sigma,\iota)$, a level $\ell<\Lambda$ and a base $X\in B_{\ell+1}$ with parent $x\in B_\ell$. Let
$\mathcal D_\ell(Y)$ denote the level-$\ell$ panels of the bases in $Y$ together with the layout bit of substep $\ell$.
\begin{enumerate}
\item[(a)] If substep $\ell$ is shrinking, then (i) the incidence list of $X$ without slot numbers, the widths of all actual arcs
  at all positions $(X,g)$, and the kernel descriptions of all positions $(X,g)$ are functions of $\mathcal D_\ell(\Ball_\ell(x,3))$;
  (ii) the full base record (P1) of $X$ is a function of $\mathcal D_\ell(\Ball_\ell(x,6))$, whose base records contain the
  addresses of all bases of $\Ball_\ell(x,7)$.
\item[(b)] If substep $\ell$ is a refresh, all of these are functions of $\mathcal D_\ell(\Ball_\ell(x,1))$.
\item[(c)] In both cases the functions are computed by a multitape machine in $\lambda^{O(1)}$ steps from the listed panels written
  on a work tape. A ball of radius $\varrho$ contains at most $(2d_\ell)^{\varrho+1}=O(1)$ bases.
\end{enumerate}
Here the \emph{kernel description} of a new position is the list of its immediate old positions, its local maps with their term
data and roles, and the slices, ties and summed fields of \eqref{eq:kernel} (H10, H12).
\end{lemma}

\begin{proof}
\emph{Part (a).} Write $d=d_\ell$ and split the level-$\ell$ lane as $(U,\eta,j,g)$ (H9(ii); we write $\eta$ for the slot
coordinate). For a base $y$ and slot $\eta$ let $f_\eta(y)$ be the outgoing base arc of slot $\eta$ at $y$ and $y^+_\eta$ its head;
for a base arc $f$ write $\tau_f=(\tau_f^U,\tau_f^\eta,\tau_f^{\rm tag},\tau_f^G)$. All of $c(y),\phi_y,f_\eta(y)$, the address of
$y^+_\eta$ and $\tau_{f_\eta(y)}$ are in the panel of $y$ (P1, P3).

\emph{Step 1 (where an old position goes).} Let $u=(y,U,\eta,j,g)$ be a level-$\ell$ position. By H9(iii),(iv) and H10:
if $U_{c(y)}=1$ (thick), $\pi(u)=((y,j),g)$. If $u$ is thin, its bundle $\beta(u)$ is the bundle of the base arc
$f=f_\eta(y)$ with source tag $j+\phi_y$ and remaining lane $g$; its packed ends are
\[
 A(\beta)=\bigl((y,\,j+\phi_y),\,g\bigr),\qquad B(\beta)=\bigl((y^+_\eta,\,j+\phi_y+\tau_f^{\rm tag}),\,g+\tau_f^G\bigr),
\]
because the thick--thick wires of $f$ from source tag $j'$ and lane $g$ end at thick positions with tag $j'+\tau_f^{\rm tag}$ and
lane $g+\tau_f^G$ at $y^+_\eta$, which are packed at $((y^+_\eta,j'+\tau^{\rm tag}_f),g+\tau_f^G)$. Then $\pi(u)=A(\beta)$ in
layout~1 and $B(\beta)$ in layout~2. Hence the parent of the base of $\pi(u)$ is $y$ or an out-neighbour of $y$, and $\pi(u)$ is
computed from the panel of $y$ and the layout bit.

\emph{Step 2 (which motifs meet a new base).} A level-$\ell$ actual arc $w$ over a base arc $e$ joins positions at bases $y_1,y_2$
($y_1\ne y_2$). From $c(y_1),c(y_2)$, the $U$-coordinates of its two ends and $\tau_e$ one reads off whether $w$ is thick--thick
(then it lies in a bundle) or not; if not, which end is $C$ (a thin end; for thin--thin, the source) and which is $D$ (H9(iii)).
Let $\mu$ be the motif of a non-bundle wire, with \emph{anchor} $a=$ base of $C$ and $y_D=$ base of $D$, so
$\mathrm{dist}_\ell(a,y_D)=1$. By Step 1 and H10 the six parties of $\mu$ are placed as follows: $A,B$ are the packed ends of
$\beta(C)$, with parents $a$ and an out-neighbour $a^+$ of $a$; $D^*=\pi(D)$ has parent $y_D$ (if $D$ is thick, or thin in
layout~1) or an out-neighbour $y_D^+$ of $y_D$ (thin, layout~2); $C$ goes to $A$ or $B$; $M,N$ go to $A$ or $B$ and to $D^*$.
So every new position touched by $\mu$ has its parent in
\[
 S(\mu)\subseteq\{a,\ a^+,\ y_D,\ y_D^+\}\subseteq\Ball_\ell(a,2),\qquad \operatorname{diam}S(\mu)\le3,
\]
the diameter bound by the path $a^+$--$a$--$y_D$--$y_D^+$. Consequently: \textbf{(2a)} if $\mu$ touches a child of $x$ then
$\mathrm{dist}_\ell(a,x)\le2$; \textbf{(2b)} the new positions of all parties of $\mu$ are computed from the panels of $a$ and
$y_D$, i.e.\ from $\mathcal D_\ell(\Ball_\ell(a,1))$.

\emph{Step 3 (the demand widths of a motif).} Let $\beta$ be the bundle of arc $f=f_\eta(a)$ with source tag $j'$ and lane $g$.
Its payload is the concatenation of the $K/4$ thick--thick wires over $f$ from the positions $(a,U',\eta',j',g)$ with
$U'_{c(a)}=1$ whose heads are thick; their widths are in the panel of $a$ (P2) and their thickness is read from $c(a^+)$ and
$\tau_f$. Its $Q$ demands $CD$ are the non-bundle wires whose $C$ is a thin position $(a,U,\eta,j,g)$ with $j+\phi_a=j'$; their
widths are in the panel of $a$ (outgoing wires) or of the in-neighbours of $a$ (incoming wires, P2 of the source). Both lists are
ordered by the identifiers of H12, written relative to the source lane of $\beta$; the payload is cut into $Q$ pieces as evenly as
possible by the fixed rule of H10, and the $c$-th piece is the $AB$ demand of the $c$-th $CD$ demand. So both demand widths of every
motif anchored at $a$, and the slices of the payload that its $AB$ demand uses, are functions of $\mathcal D_\ell(\Ball_\ell(a,1))$.

\emph{Step 4 (terms).} The term index of a motif anchored at $a$ is in the panel of $a$ (P3). Its admission and its supply widths
$c_{\alpha,e}$ follow from the index and the two demand widths by H6 in $\mathrm{poly}(\log\lambda)$ bit operations.

\emph{Step 5 (surviving arcs).} Given the layout bit, the surviving group is $G_2$ (layout~1) or $G_1$ (layout~2). Each surviving
supply type $e$ of $\mu$ gives $m_e$ parallel actual arcs between the new positions of its two roles (Step 2), oriented by the
fixed role order, with widths the even split of $c_{\alpha,e}$ (H10). Its identifier is (the $CD$ arc of $\mu$ with lanes relative
to the endpoint being numbered, $e$, the subdivision index) (H9(vii)). All of this is a function of
$\mathcal D_\ell(\Ball_\ell(a,1))$.

\emph{Step 6 (dummy arcs).} By H9(vi) and H10 every new position has the same deficit $d'-dQ\sum_em_e$, a constant. The dummy arcs
at $((y,j),g)$ go to $((y,j^{\rm next}),g)$ and come from $((y,j^{\rm prev}),g)$ for the cyclic order of tags (H9(vii)), which we take to be $j^{\rm next}=j+1$ and $j^{\rm prev}=j-1$ modulo $J$ in the binary
numbering of (R2).
They need no data.

\emph{Step 7 (proof of (a)(i)).} The positions $(X,g)$, $X=(x,j)$, receive: old positions $u$ with $\pi(u)=(X,g)$, which by
Step~1 lie at bases in $\Ball_\ell(x,1)$; roles of motifs with anchor $a\in\Ball_\ell(x,2)$ (Step 2a), each computed from
$\mathcal D_\ell(\Ball_\ell(a,1))\subseteq\mathcal D_\ell(\Ball_\ell(x,3))$ (Steps 2b--5); and dummy arcs (Step 6). Enumerate all
positions of all bases of $\Ball_\ell(x,2)$ and all motifs anchored there, compute their placements, and keep those touching a
position above $X$. This gives the unnumbered incidence list of every $(X,g)$ with its identifiers and widths (Step 5). The kernel
description of $(X,g)$ consists of the old positions placed there (Step 1), the maps of the roles placed there with their term
data (Steps 2--4), the slices of bundle payloads (Step 3) and of surviving payloads (Step 5), the ties of the supply edges of the
contracted (non-retained) group, which became local, and the identification of each placed old position's half-wires with the demand fields of its
motifs. All are among the data just computed. This proves (a)(i).

\emph{Step 8 (proof of (a)(ii)).} By H9(vii) and H12 the slot numbers at $X$ are the ranks of the identifiers in the incidence
list of $(X,0)$, sorted separately for incoming and outgoing arcs; translation of the remaining lane preserves them, so these are
the slot numbers at every lane. Consider an arc at $(X,0)$ produced by a motif $\mu$ (or a dummy), with other end $(X',g')$,
$X'=(x',j')$. Both $x$ and $x'$ lie in $S(\mu)$ (or $x'=x$ for a dummy), so $\mathrm{dist}_\ell(x,x')\le3$. The slot number of the
arc at $X'$ is its rank in the incidence list of $(X',g')$, which equals the rank of the translated arc in the list of $(X',0)$;
by Step 7 applied to $X'$ this list is a function of $\mathcal D_\ell(\Ball_\ell(x',3))\subseteq\mathcal D_\ell(\Ball_\ell(x,6))$.
The displacement of the base arc is $g'$ (an arc from $(X,0)$ ends at lane $g'$; an arc into $(X,0)$ from $(X',g')$ has
displacement $g'$ as well, since $g'+g'=0$ in $\F_2^h$). The address of $X'$ is $\mathrm{addr}(x')J_\ell+j'$, and the address of
$x'\in\Ball_\ell(x,3)$ is in the base records of $\Ball_\ell(x,2)$. This proves (a)(ii). (Radius $6$ suffices; the base records of
$\Ball_\ell(x,6)$ name the bases of $\Ball_\ell(x,7)$, which are never needed beyond their addresses.)

\emph{Part (b).} Split the level-$\ell$ lane as (port $j$, remaining $g$) (H11). For an old base arc $e:x\to y$ with outgoing slot
$o(e)$ at $x$ and incoming slot $i(e)$ at $y$ (the latter stored in the base record of $x$), the external arc of $e$ at remaining
lane $g$ runs from $((x,o(e)),g)$ to $((y,i(e)),g+\tau^G_e)$ and its width is the sum of the $D$ widths of the old arcs over $e$ from
the block $(x,g)$ (panel of $x$). The external arc entering port $j$ of $(x,g)$ is the external arc of the old arc $e'$ in incoming
slot $j$ of $x$, whose source $y'$ and outgoing slot $o(e')$ are in the base record of $x$ and whose width is in the panel of
$y'\in\Ball_\ell(x,1)$. Internal arcs join $(x,j)$ to $(x,j')$ along the fixed port graph, displacement $0$, widths given by the block
term (panel of $x$) and H6. Padding arcs go to the next port cyclically. The external slot is slot $0$ in both directions, so the far
slot number of an external arc is $0$. The kernel of $((x,j),g)$ holds the single old position $(x,(j,g))$, the ties of its old
half-wires to their segments of the external wires (the offsets are prefix sums of the old widths in source-port order, from the
panels of $x$ and of the in-neighbours of $x$), and the internal maps of the block term. Everything is a function of
$\mathcal D_\ell(\Ball_\ell(x,1))$.

\emph{Part (c).} The undirected degree of $B_\ell$ is at most $2d_\ell$, so a ball of radius $\varrho$ has at most
$(2d_\ell)^{\varrho+1}$ bases, a constant. Each panel has $\lambda^{O(1)}$ bits. The computations above enumerate at most
$O(1)\cdot N=\lambda^{O(1)}$ positions and motifs, sort bounded lists of identifiers of $O(\lambda)$ bits, add and compare
integers of $O(\lambda)$ bits, and evaluate H6's admission test; on a multitape machine each of these costs $\lambda^{O(1)}$.
\end{proof}

\begin{remark}
The radii in Lemma~\ref{lem:locality} improve the radius $10$ asserted in the first written version of the port case and agree with the
bounds $3$ and $6$ found in its line-by-line check. Any fixed radius would do: only $O(1)$ batched lookup rounds per level are needed.
\end{remark}

\subsection{The witness and its checks}\label{subsec:witness}

\paragraph{Fixed choices.} Where EG leaves a choice free we fix one deterministic rule. For the colours and tag offsets we use
exactly the greedy scans of OpenAI's Lean construction (\texttt{orderedColor} and \texttt{orderedTagOffsets},
\src{OAI/Computability/VPVNP/SlotGreedyScan.lean} l.182 and l.235): bases are visited in increasing address, starting from tables
that hold $0$ (colour $0$, tag $0$) at every base, and each visit reads the current tables at the $2d$ neighbour slots of $y$. So a
neighbour of \emph{larger} address is read as having value $0$. Write $N(y)$ for the bases at the slots of $y$, and for
$y'\in N(y)$ put $\tilde c_y(y')=c(y')$, $\tilde\phi_y(y')=\phi_{y'}$ if $\mathrm{addr}(y')<\mathrm{addr}(y)$, and
$\tilde c_y(y')=0$, $\tilde\phi_y(y')=0$ otherwise.
\begin{itemize}
\item[(R1)] $c(y)$ is the least colour in $\{0,\dots,2d\}$ outside $\{\tilde c_y(y'):y'\in N(y)\}$ (at most $2d$ values, so it exists).
\item[(R2)] $\phi_y$ is the first tag, in increasing order of the binary number that a tag (a vector in $\F_2^{\log J}$) spells
  with its first coordinate most significant (the numbering \texttt{bitWordFinEquiv} of OpenAI's Lean; the proof needs only that the
  construction and the checks use the same order), outside the list of at most
  $2d((d+1)+(d+1)^2)<J$ forbidden values: at each incidence of a base arc $e$ at $y$, with other end $y'$, the values
  $\tau^{\rm tag}_e+\ell$ ($\ell\in L_y$) and $\tau^{\rm tag}_e+\ell+\ell'+\tilde\phi_y(y')$ ($\ell\in L_y$, $\ell'\in L_{y'}$).
\end{itemize}
These scans give a proper colouring and offsets that satisfy H9(v) at every incidence: the condition between $y$ and an earlier
neighbour $y'$ is imposed when $y$ is visited, with the final $c(y')$, $\phi_{y'}$; reading later neighbours as $0$ only forbids
additional values. Payloads are cut with the larger pieces first; the layout tie goes to layout~1; the external slot of a refresh
is slot $0$, internal slots follow the port-graph order, dummies come last. By the scope stated in H8, which rests on EG's
proof, H8 applies to these choices. The four rules above (the two greedy scans, the larger-pieces-first split, the tie to layout~1)
are also the rules of OpenAI's Lean construction (\texttt{orderedColor}, \texttt{orderedTagOffsets}, \texttt{splitWidth},
\texttt{chosenLayout}); we do not claim that the rest of our construction matches the Lean one (Remark~\ref{rem:weak}, item 3).

\paragraph{Energy scores.} For a shrinking substep and a motif $\mu$ with term $\alpha$, let $e_k(\mu)$ ($k=1,2$) be the energy
of the group $G_k$ of $\mu$, written after clearing the fixed denominator as an integer combination, with fixed coefficients, of the
squares of the supply widths $c_{\alpha,e}$, $e\in G_k$ (H12). For a base $y$ put $E_k(y)=\sum_{y'\le y}\sum_{\mu\text{ anchored
at }y'}e_k(\mu)$. Layout~1 (which retains $G_2$) is chosen iff $E_2(y_{\max})\le E_1(y_{\max})$, by the fixed tie rule.

\paragraph{The witness.} A witness is a string $u\in\{0,1\}^m$ with the following sections at fixed offsets:
\begin{enumerate}
\item[S0] $\sigma$, two bits per trit ($00\mapsto0$, $01\mapsto1$, $10\mapsto-1$; $11$ is invalid), $2rIN$ bits;
\item[S1] the \emph{directory}: for each $\iota,\rho$ the start block $\beta_\rho(\iota N)$, and for each $\delta$ the previous and
  the next owner of the big block $q_\rho(\iota)+\delta$ (or ``none''), each in $\lambda+O(1)$ bits;
\item[S2] for $\ell=0,\dots,\Lambda$: the layout bit of substep $\ell$ (if shrinking) and the panels of all bases of $B_\ell$, in
  address order, each of size $\Pi_\ell$;
\item[S3] the string $z$, followed by zeros up to length $m$.
\end{enumerate}
The sizes of $B_\ell$ and $\Pi_\ell$ depend only on $n$ and the fixed data, so the offsets of S0--S3 and of every panel are
computable from $n$. Given $u$, write $\sigma(u)$ for the content of S0 and $\iota(u)$ for the term indices in S2.

\paragraph{The checks.} The verifier $V(x,u,\gamma)$ reads $\gamma$ as a number and performs the check with that number (in a
fixed enumeration of the checks below; numbers beyond the last check accept). Each check also verifies that the fields it reads
are in range and that their unused parts are zero.
\begin{enumerate}
\item[C1] (one per trit) The trit code is not $11$.
\item[C2] (one per $\rho$, and one per $\rho$ and $\iota<I-1$) $\beta_\rho(0)=0$, and
  $\beta_\rho((\iota+1)N)=\beta_\rho(\iota N)+\sum_{g<N}\sigma_\rho(\iota N+g)$ (the value $\beta_\rho(IN)$ is not stored).
\item[C3] (one per $\iota,\rho,\delta$) With $Q=q_\rho(\iota)+\delta$, the stored next owner is the least $\iota'>\iota$ with
  $Q\in q_\rho(\iota')+\{-1,0,1\}$, or ``none'' if there is none; likewise for the previous owner.
\item[C4] (one per level-$0$ base) Its panel (P1), (P2) equals the one computed from the directory entries of its big interval
  and of the at most $6r$ owners named in them (these give the far ends of the owner links),
  the piece of $\sigma$ of that interval, and the construction of Section~\ref{subsec:network} (including the looping algorithm).
\item[C5] (one per base of a shrinking level) $c(y)$ equals the value of (R1) computed from the stored colours of the bases of
  $N(y)$ of smaller address, with $0$ for those of larger address (addresses from the base record of $y$).
\item[C6] (one per base of a shrinking level) $\phi_y$ equals the value of (R2) computed in the same way from the stored offsets
  and the displacements in the base records of $y$ and of the bases of $N(y)$.
\item[C7] (one per term slot) The term slot is zero if no motif (block) is anchored there; otherwise it holds the fixed-length
  encoding of an index admitted (H6) for the demand widths computed by Steps 2--3 of Lemma~\ref{lem:locality} (resp.\ by the block
  widths, for a refresh).
\item[C8] (one per base of a shrinking level) $E_k(y)=E_k(y-1)+\sum_{\mu\text{ anchored at }y}e_k(\mu)$, with $E_k(-1)=0$.
\item[C9] (one per shrinking level) The layout bit agrees with $E_1,E_2$ at the last base and the tie rule.
\item[C10] (one per base $X$ of each level $\ell+1\le\Lambda$) The panel entries (P1), (P2) of $X$ equal the values computed by
  Lemma~\ref{lem:locality} from the level-$\ell$ panels of $\Ball_\ell(\mathrm{parent}(X),6)$ and the layout bit.
\item[C11] (one per final base $y$) If $y$ is the first final base, $\Omega(y)=0$. If $y$ is not the last final base,
  $\Omega(y+1)=\Omega(y)+\sum_{g,\eta}(\text{width of out-slot }\eta\text{ at }(y,g))$. The last base has no successor and no
  stored $\Omega(y+1)$ (P4 stores prefixes only for actual bases); its sum is checked by C12.
\item[C12] (one) Let $L=\Omega(y_{\max})+\sum(\text{widths at }y_{\max})$. Then $L$ is at most the length of S3, and every bit of S3
  at a position $i\ge L$ is $0$ (positions are counted from $0$, so the bit at position $L$ is checked too).
\item[C13] (one per final position $v$) $W_v(z|_v;\iota(u))=1$, where $z|_v$ is read from S3 at the offsets given by (P4) and (P2)
  (the bits of an arc are stored once, at its source position, in the order base, lane, outgoing slot).
\end{enumerate}

\begin{lemma}[Uniqueness]\label{lem:unique}
Let $n\ge n_M$. For every $\sigma\in\{-1,0,1\}^{rIN}$, $\iota\in\mathrm{Adm}(\sigma)$ and $z\in\{0,1\}^{\ell(\sigma,\iota)}$ there is
exactly one witness $u=u(\sigma,\iota,z)$ that passes C1--C12 and has $\sigma(u)=\sigma$, $\iota(u)=\iota$ and $z$ at the start of
S3. Every witness that passes C1--C12 is of this form. For $u=u(\sigma,\iota,z)$, the check C13 at $v$ evaluates the true vertex
tensor $W_v(z|_v;\iota)$ of \eqref{eq:expansion}.
\end{lemma}
\begin{proof}
\emph{Existence.} Write into $u$ the encoding of $\sigma$, the true directory, and for each level the true panels of the level-$\ell$
network determined by $(\sigma,\iota)$ with the fixed choices, the true layout bits, energy prefixes and final prefixes, then $z$
and zeros. Lemma~\ref{lem:space} below shows that this fits into $m$ bits. Every check passes: C4 and C10 by construction and
Lemma~\ref{lem:locality}, the others by definition.

\emph{Uniqueness and form.} Let $u$ pass C1--C12. By C1, $\sigma(u)$ is a skeleton. By C2 and C3 (induction over $\iota$) the
directory is the true one for $\sigma(u)$; by C4 the level-$0$ panels are the true base records and widths of $N_{x,\sigma(u)}$.
Suppose the level-$\ell$ base records and widths in $u$ are the true ones for $\sigma(u)$ and the terms $\iota_{<\ell}(u)$ of the
earlier substeps. By C5 and C6, induction in address order shows that the colours and offsets are the ones of the scans (R1),
(R2): if the stored values agree with the scans at all bases of smaller address than $y$, then C5 and C6 at $y$ compute (R1) and
(R2) from exactly the table that the scan reads when it visits $y$ (final values below $\mathrm{addr}(y)$, the value $0$ above it),
so the stored values at $y$ agree as well. By C7
the terms of substep $\ell$ are admitted for the true demand widths, so $\iota_{\le\ell}(u)$ is an admitted prefix. By C8 and C9
the energy prefixes and the layout bit are the true ones. By C10 and Lemma~\ref{lem:locality} the level-$(\ell+1)$ base records and
widths are the true ones. At the end $\iota(u)\in\mathrm{Adm}(\sigma(u))$, the final prefixes are true by C11, and by C12 S3 is
$z$ followed by zeros, with $|z|=\ell(\sigma(u),\iota(u))$. So $u=u(\sigma(u),\iota(u),z)$; and $u(\sigma,\iota,z)$ is determined by
$(\sigma,\iota,z)$ because each field is the true value or a forced zero.

\emph{C13.} For such $u$ the verifier rebuilds the recursive expansion of $v$ from true panels (Lemma~\ref{lem:locality}(a)(i),(b),
level by level down to level $0$, where the initial factors of Section~\ref{subsec:network} are recomputed from $x$, the directory and
$\sigma$), and sums over the hidden indices of $v$ with $z|_v$ fixed. By H8 (last sentence) this is $W_v(z|_v;\iota)$.
\end{proof}

\begin{lemma}[Space]\label{lem:space}
For $n\ge n_M$, the sections S0--S2 together with any $z$ of length $\ell(\sigma,\iota)$ fit into $m$ bits.
\end{lemma}
\begin{proof}
S0 has $2rIN=O_M(t/B)=O_M(t\lambda^{-1/2})$ bits and S1 has $O_M(I\lambda)=O_M(t\lambda/b)=O(t\lambda^{-9})$ bits. For S2 let
$P_\ell=|B_\ell|\cdot(\#\text{lanes at level }\ell)$ be the number of positions at level $\ell$. By H8, all levels of block $i$
have at most $P_0\kappa_0^{-i}$ positions, and each block has $s_0+1$ levels, so $\sum_\ell P_\ell\le2(s_0+1)P_0$. Since every level
has at least $(\lambda/2)^9$ lanes \eqref{eq:p2}, $\sum_\ell|B_\ell|=O(P_0\lambda^{-9})$. The per-base fields (base record, colour,
offset, both energy prefixes, the final prefix: $O(d_\ell\lambda)$ bits) therefore cost $O(P_0\lambda^{-8})$ bits in all, and the
per-position fields (widths and term slots, $O(\log\lambda)$ bits each, a bounded number per position) cost $O(P_0\log\lambda)$.
By \eqref{eq:size} the total of S2 is $O_M(t(\log\lambda)^2\lambda^{-1/2})$. Since $s\le1/64<1/2$ and
$m\ge t\lambda^{-s}(\log\lambda)^6$, S0--S2 take at most $m/2$ bits for $n\ge n_M$. By \eqref{eq:zlen},
$|z|\le c_M2^{-T}t\log\lambda\le m\,c_M\log\lambda/(p+1)^6\le m/2$ for $n\ge n_M$.
\end{proof}

\begin{lemma}[Time]\label{lem:time}
For $n\ge n_M$, $V$ halts within $O_M(m\log\lambda)$ steps on every $(x,u,\gamma)$ with $|x|=n$, $|u|=m$, $|\gamma|=k$.
\end{lemma}
\begin{proof}
Computing $t,\lambda,p,B,h,N,b,T,I,m$, the section offsets and $\Pi_\ell$ costs $O(n)+\lambda^{O(1)}$ steps (count $n$ by
Lemma~\ref{lem:countdown}, then binary arithmetic). Every fetch below is a batched lookup (Lemma~\ref{lem:batch}) of
$\lambda^{O(1)}$ requests of total length $\lambda^{O(1)}$, costing $O(m)+\lambda^{O(1)}$.

C1, C2, C5--C9, C11 use $O(1)$ rounds and $\lambda^{O(1)}$ computation. C3 scans the directory between $\iota$ and the stored owner
and computes $\lfloor\beta/N\rfloor$ (a shift by $h$ bits) per entry: $O(I\lambda)=o(m)$. C4 fetches the piece of $\sigma$ of one big
interval ($rN$ trits), the directory entries of that interval and of the at most $6r$ owners named there, one panel, and recomputes the access permutation, the Bene\v{s} settings and the
regularization in $\mathrm{poly}(N)=\lambda^{O(1)}$. C10 fetches $\Ball_\ell(\mathrm{parent}(X),6)$ by $7$ rounds of breadth-first
search over the addresses in the base records, then applies Lemma~\ref{lem:locality}(c). C12 is one scan.

C13: starting from $v$ at level $\Lambda$, for $\ell=\Lambda-1,\dots,0$ the verifier fetches, in $7$ rounds, the level-$\ell$ balls of
radius $6$ around the parents of the occurrences found so far, and computes their kernel descriptions (Lemma~\ref{lem:locality}). By
H8 there are at most $C_2^T\le\lambda^{1/64}$ occurrences in the whole expansion, so each round has $\lambda^{O(1)}$ requests. Then
it fetches $z|_v$ (one round, offsets from (P4) and (P2) of the bases that carry the arcs of $v$), and for the level-$0$ factors
the relevant pieces of $\sigma$, directory entries and input bits (one scan of $x$, $O(n)$). If the description claims more than
$\lambda^{3/4}$ hidden bits, or more than $C_2^T$ occurrences, it rejects. Otherwise it enumerates all assignments of the at most
$\lambda^{3/4}$ claimed hidden bits (for a true expansion there are at most $c_M\lambda^{17/32}$ of them, by \eqref{eq:hidden}, so
true expansions are never rejected for large $n$); for each it evaluates at most $C_2^T$ occurrences: map entries by H6
($\lambda^{O(1)}$, with $\kappa\in\F_2$ hardwired, so a monomial is $0$ or $1$), ties, pins and routing factors, and interval tables
$\Phi_i$ by enumerating $2^{e_MB}$ code strings. This costs $2^{O_M(\lambda^{3/4})}=t^{o(1)}$ steps. In all C13 costs
$O((\Lambda+1)(m+n))+t^{o(1)}=O_M(m\log\lambda)$, since $\Lambda=(s_0+1)T=O(\log\lambda)$ and $n\le m$.
\end{proof}

\begin{proof}[Proof of Lemma~\ref{lem:P}]
Let $n\ge n_M$ and $x\in\{0,1\}^n$. The number of checks is $O(\sum_\ell P_\ell+rIN)=O_M(t\log\lambda/B)\le t$, so $k=\lceil\log
t\rceil$ bits address them all. By Lemma~\ref{lem:unique}, a witness passes every check iff it equals $u(\sigma,\iota,z)$ for some
$\sigma$, $\iota\in\mathrm{Adm}(\sigma)$, $z$, with $W_v(z|_v;\iota)=1$ for all final $v$. Since $W_v\in\{0,1\}$,
\[
 \#\{u:\forall\gamma\ V(x,u,\gamma)=1\}=\#\{(\sigma,\iota,z):\textstyle\prod_vW_v=1\}
 \equiv\sum_\sigma\sum_{\iota\in\mathrm{Adm}(\sigma)}\sum_z\prod_vW_v(z|_v;\iota)\pmod2,
\]
which by \eqref{eq:expansion} and Lemma~\ref{lem:pathcount} is $\sum_\sigma Z_\sigma(x)\equiv\acc(M,x)\pmod2$. By
Lemma~\ref{lem:time}, $V$ runs in time $O(m\log\lambda)$.
\end{proof}

%======================================================================
\section{Proof of the main theorem}\label{sec:wrapper}

We write $\NTIME[n;2\log n]$ for $\NTIME[f;g]$ with $f(n)=n$, $g(n)=\lfloor2\log n\rfloor$. If $g\le g'=O(f)$ and a machine that reads
its input can compute $g(n)$ within $O(f(n))$ steps (true for all budgets here), then $\NTIME[f;g]\subseteq\NTIME[f;g']$: the new verifier copies the first
$g(n)$ guess bits to a fresh tape, followed by a blank, and runs the old one.

\begin{proposition}[Case 2]\label{prop:case2}
If characteristic $2$ has ports, then
\[\NTIME[n;2\log n]\not\subseteq\DTIME(n(\log n)^{s/2}),\]
with $s$ as in Section~\ref{sec:case2}.
\end{proposition}
\begin{proof}
Suppose $\NTIME[n;2\log n]\subseteq\DTIME(n(\log n)^{s/2})$. Let $A\in\parTIME(n^2)$, and take $V,m,k$ from Lemma~\ref{lem:P}, with
$\lambda,p$ as there. Put $R(n)=(p+1)\,m(n)$ and
\[
 A''=\bigl\{\,1^n0\,x\,u\,0^{R(n)}:\ n\ge n_M,\ x\in\{0,1\}^n,\ u\in\{0,1\}^{m(n)},\ \forall\gamma\in\{0,1\}^{k(n)}\
 V(x,u,\gamma)=1\,\bigr\}.
\]
Let $w$ have length $\nu$. \emph{The complement of $A''$ is in $\NTIME[\nu;2\log\nu]$:} count the leading ones to get $n$; if $w$ is
not of the form $1^n0xu0^{R(n)}$ with $n\ge n_M$ (a check of lengths and of the zero tail, using counters, $O(\nu)$ steps), accept;
otherwise take the first $k(n)$ guess bits as $\gamma$ and accept iff $V(x,u,\gamma)=0$. Here $\nu\ge m\ge t^{0.9}$, so
$k=\lceil\log t\rceil\le\lfloor2\log\nu\rfloor$ for large $n$, and the time is $O(\nu)+O(m\log\lambda)=O(\nu)$ because
$p+1\ge\log\lambda$. By the assumption and closure of $\DTIME$ under complement, some deterministic machine $E$ decides $A''$ in time
$O(\nu(\log\nu)^{s/2})$.

\emph{A faster parity machine for $A$.} On input $x$ with $n\ge n_M$: write $1^n0x$, write $m(n)$ nondeterministic bits $u$, write
$0^{R(n)}$, run $E$ on the resulting string and accept iff $E$ accepts (for $n<n_M$ use a finite table, with one accepting
computation iff $x\in A$). Its accepting computations correspond to the $u$ with $1^n0xu0^{R}\in A''$, so their number is
$\#\{u:\forall\gamma\,V(x,u,\gamma)=1\}$, which is odd iff $x\in A$ by Lemma~\ref{lem:P}. Its running time is
$O(\nu(\log\nu)^{s/2})$ with $\nu=2n+1+(p+2)m=O(t\lambda^{-s}(\log\lambda)^7)$ and $\log\nu\le\lambda+O(1)$, i.e.
\[
 O\bigl(t\,\lambda^{-s/2}(\log\lambda)^7\bigr)=o(n^2).
\]
So $\parTIME(n^2)\subseteq\parTIME(t'')$ with $t''(n)=n^2(\log n)^{-s/2}(\log\log n)^7=o(n^2)$, which contradicts
Theorem~\ref{thm:hier} with $t(n)=n^2$.
\end{proof}

\begin{proof}[Proof of Theorem~\ref{thm:main}]
If characteristic $2$ has no ports, Corollary~\ref{cor:case1} gives $\NTIME[n;\lfloor\log n\rfloor+1]\not\subseteq\DTIME(n(\log n)^{1/4})$, and
$\NTIME[n;\lfloor\log n\rfloor+1]\subseteq\NTIME[n;2\log n]$, since $\lfloor\log n\rfloor+1\le\lfloor2\log n\rfloor$ for $n\ge2$. If it has ports, Proposition~\ref{prop:case2} applies. So the theorem
holds with $\sigma=1/4$ or $\sigma=s/2$, according to the case. Finally $\NTIME[n;2\log n]\subseteq\NTIME(n)$.
\end{proof}

%======================================================================
\section{Nondeterministic and unambiguous machines}\label{sec:utime}

The two cases of the proof extend from deterministic machines to two larger classes, one in each case: Case~1 to all
nondeterministic machines, Case~2 to parity machines. Since $\mathrm{UTIME}(f)\subseteq\NTIME(f)\cap\parTIME(f)$, this gives a
separation from unambiguous time. The classes are those of Definition~\ref{def:classes}.

\begin{theorem}\label{thm:Uplus}
There are $\sigma>0$ (not effective) and a language $L\in\coNTIME[n;2\log n]$ such that
$L\notin\NTIME(n(\log n)^{\sigma})\cap\parTIME(n(\log n)^{\sigma})$. In particular
$\coNTIME[n;2\log n]\not\subseteq\mathrm{UTIME}(n(\log n)^{\sigma})$.
\end{theorem}

As before, $\sigma=1/4$ if characteristic $2$ has no ports and $\sigma=s/2$ if it has ports. The language $L$ depends on the case:
$L=\TRANS$ in the first; in the second, the padded language $A''$ of Lemma~\ref{lem:B}, built from a language $A$ that witnesses
the parity time hierarchy at time $n^2$. Theorem~\ref{thm:Uplus} implies
Theorem~\ref{thm:main}: if $\NTIME[n;2\log n]\subseteq\DTIME(n(\log n)^{\sigma})$, then, $\DTIME$ being closed under complement,
$\coNTIME[n;2\log n]\subseteq\DTIME(n(\log n)^{\sigma})\subseteq\mathrm{UTIME}(n(\log n)^{\sigma})$.

\begin{lemma}[Lemma A]\label{lem:A}
Assume characteristic $2$ has no ports. Let $M$ be an NTM such that every input $(X,X^T)$ has an accepting computation and no input
$(X,Y)$ with $Y\ne X^T$ has one ($X,Y\in\{0,1\}^{N\times N}$, row-major, $N\ge1$). Then there are $c_M>0$ and $N_M$ such that for
every $N\ge N_M$ some $X$ has no accepting computation on $(X,X^T)$ with fewer than $c_MN^2\sqrt{\log N}$ steps.
\end{lemma}

\begin{proof}
Fix $N$ and let $T$ be the largest, over $X$, of the length of a shortest accepting computation on $(X,X^T)$; assume
$T\le N^2\sqrt{\log N}$ (otherwise there is nothing to prove). Give the transitions of $M$ codes as in
Section~\ref{subsec:normal}: exactly one valid code in $\{0,1\}^{e_M}$ for each available transition, and one code that keeps a
halted configuration unchanged. Let $E=\max(1,\lceil T/N\rceil)$. Then each computation of length at most $EN$ extends to exactly one
\emph{run}, a sequence of $EN$ valid codes, and the run ends in an accepting state iff the computation accepts.

\emph{Records.} Blocks of $N$ cells and epochs of $N$ steps are as in the proof of Proposition~\ref{prop:trans}. For every $X$ fix
one accepting computation $\pi_X$ on $(X,X^T)$ of length at most $T$. The record of a run is the list of the blocks scanned by each
head at each epoch boundary; there are at most $3^{rE}$ records ($r$ tapes). The choice codes are not part of the record. Let $R$
be the record of the largest number of the runs of the $\pi_X$, and $\mathcal G_R$ the set of those $X$:
$|\mathcal G_R|\ge3^{-rE}2^{N^2}$.

\emph{Network with formal weights.} Take the graph, the terminals, the wires and all vertices other than the epoch vertices exactly as
in the proof of Proposition~\ref{prop:trans}, for the record $R$. For each epoch $j$ and each word $w\in(\{0,1\}^{e_M})^N$ let
$z_{j,w}$ be an indeterminate. The table of epoch $j$ at given incident strings is the sum of $z_{j,w}$ over all words $w$ that are
valid codes at every step and lead in $N$ steps from the incoming data (state, head offsets, contents of the blocks of the window)
to the outgoing data, with the block changes prescribed by $R$; for the first epoch the incoming state and head offsets are the
initial ones, and for the last epoch the words must end in an accepting halted configuration. For external strings $X$ at the $A_i$ and $Y$ at the $B_k$ let $P_{X,Y}$ be the contraction, a polynomial over
$\F_2$ in the $z_{j,w}$. Expanding the tables, $P_{X,Y}$ is the sum, over pairs (labelling of the internal strings, words
$w_1,\dots,w_E$) at which every non-epoch table is $1$ and each $w_j$ is admitted by its table, of $\prod_jz_{j,w_j}$. Exactly as in
the proof of Proposition~\ref{prop:trans} (induction over epochs in time order; the contents of a block arrive from its last earlier
use, and between two uses no head is in the block), these pairs are in bijection with the runs on $(X,Y)$ that follow $R$ and
accept: a run determines its words and every label (summed final contents included, so each run is counted once), and the words of
a pair concatenate to such a run. On a fixed input the code word sequence determines the run, so distinct runs give distinct
monomials. Hence
\[
 P_{X,Y}=\sum_{\text{runs on }(X,Y)\text{ that follow }R\text{ and accept}}\ \prod_{j=1}^{E}z_{j,w_j(\text{run})},
\]
a sum of distinct multilinear monomials of degree $E$ with coefficient $1$. If $Y\ne X^T$ there is no accepting computation, so
$P_{X,Y}=0$ identically. If $X\in\mathcal G_R$, the run of $\pi_X$ contributes, so $P_{X,X^T}\ne0$.

\emph{Specialization over $\F_2$.} Let $X\in\mathcal G_R$ and let $\mu$ be a monomial of $P_{X,X^T}$. Fix the indeterminates that
do not occur in $\mu$ to arbitrary values in $\F_2$. Every other monomial of $P_{X,X^T}$ misses some indeterminate of $\mu$ (it is
multilinear, of degree $E$ and different from $\mu$), so what remains is a multilinear polynomial in the $E$ indeterminates of $\mu$
in which $\mu$ has coefficient $1$. Its sum over $\{0,1\}^E$ equals that coefficient, because the sum of a monomial over the cube is
even unless the monomial contains every indeterminate. So the polynomial takes the value $1$ somewhere. Hence, for $z$ uniform over
all assignments in $\F_2$, $\Pr[P_{X,X^T}(z)=1]\ge2^{-E}$, and by averaging one assignment $z$ satisfies $P_{X,X^T}(z)=1$ for all
$X$ in a set $\mathcal G\subseteq\mathcal G_R$ with $|\mathcal G|\ge2^{-E}|\mathcal G_R|\ge2^{N^2-(1+r\log3)E}$. Fix this $z$.
The tables are now $\F_2$-valued functions of the incident strings.

\emph{Conclusion.} For $X\in\mathcal G$, fixing the $A_i$ to the rows of $X$ makes the contraction, as a tensor in the column
indices, equal to $\bigotimes_ke_{\mathrm{col}_k(X)}$, because $P_{X,Y}(z)$ is $1$ at $Y=X^T$ and $0$ at every other $Y$. Since
$E=O(1+T/N)=O(N\sqrt{\log N})=o(N^2)$, $|\mathcal G|\ge2^{(1-\rho)N^2}$ with $\rho=1/4$ for $N\ge N_M$. The graph is that of
Proposition~\ref{prop:trans}, so it has maximum degree $O_M(1)$ in bundles, bundle widths $O_M(N)$ and $O_M(N+1+T/N)$ vertices.
Proposition~\ref{prop:table} gives $O_M(N+1+T/N)\ge cN\sqrt{\log N}$, so $T\ge c'_MN^2\sqrt{\log N}$ for large $N$.
\end{proof}

\begin{remark}
(1) We index the weights by an epoch and a whole code word. With one weight per transition the product would be commutative,
two runs with the same multiset of transitions would share a monomial, and its coefficient could be even. (2) Over an infinite field
one could keep all of $\mathcal G_R$ by choosing $z$ outside the zero set of $\prod_{X\in\mathcal G_R}P_{X,X^T}$; this is the idea of
OpenAI's Lean lemma \texttt{realize\_relations} (one indeterminate per vertex and local view, over an infinite field). But then Proposition~\ref{prop:table} is needed over that field with
nonzero scalars $c(X)$ in place of $1$, that is, H2, H4 and H5 for that field. (For the algebraic closure of $\F_2$ the class of spaces is $\X_2$ again, by H1(b) and
H2, so H4 would apply as stated; the price comparison and the entropy lemma with general nonzero scalars would be new inputs.) The specialization over $\F_2$ avoids this, so
the proof of Theorem~\ref{thm:Uplus} uses the inputs H1--H12 exactly as stated. (3) For an unambiguous machine one may take every $z_{j,w}=1$: the
contraction is then the number of accepting runs with record $R$, which is $0$ or $1$.
\end{remark}

\begin{corollary}\label{cor:A}
If characteristic $2$ has no ports, then $\TRANS\in\coNTIME[n;\lfloor\log n\rfloor+1]$ and $\TRANS\notin\NTIME(n(\log n)^{\sigma})$ for every
$\sigma<1/2$.
\end{corollary}
\begin{proof}
The complement of $\TRANS$ is in $\NTIME[n;\lfloor\log n\rfloor+1]$ (Corollary~\ref{cor:case1}). An NTM for $\TRANS$ that runs in time
$O(n(\log n)^{\sigma})$ has, at $n=2N^2$, only computations of length $O(N^2(\log N)^{\sigma})=o(N^2\sqrt{\log N})$, which
contradicts Lemma~\ref{lem:A}.
\end{proof}

\begin{lemma}[Lemma B]\label{lem:B}
Assume characteristic $2$ has ports, and let $s$ be as in Section~\ref{sec:case2}. Let $A$ be a language given by
Theorem~\ref{thm:hier} for $t(n)=n^2$ (so $A\in\parTIME(n^2)$ and $A\notin\parTIME(t')$ for every $t'=o(n^2)$), and define
$V,m,k,R(n)$ and $A''$ for this $A$ as in the proof of Proposition~\ref{prop:case2}. Then $A''\in\coNTIME[\nu;2\log\nu]$ and
$A''\notin\parTIME(\nu(\log\nu)^{s/2})$.
\end{lemma}
\begin{proof}
The proof of Proposition~\ref{prop:case2} shows that the complement of $A''$ is in $\NTIME[\nu;2\log\nu]$. Suppose some NTM $P$
runs in time $O(\nu(\log\nu)^{s/2})$ with $w\in A''$ iff $\acc(P,w)$ is odd. The parity machine $P_A$, on input $x$ with $n\ge n_M$, writes $1^n0x$, then $m(n)$ nondeterministic bits $u$ with one binary
branching per bit (so each $u$ has exactly one computation prefix), then $0^{R(n)}$, and runs $P$ on $w_u=1^n0xu0^{R(n)}$, accepting iff
$P$ accepts; for $n<n_M$ it uses a finite table, with one accepting computation iff $x\in A$. Its accepting computations are the pairs
($u$, accepting computation of $P$ on $w_u$), so
\[
 \acc(P_A,x)=\sum_u\acc(P,w_u)\equiv\#\{u:\ w_u\in A''\}=\#\{u:\ \forall\gamma\ V(x,u,\gamma)=1\}\pmod2,
\]
which is odd iff $x\in A$ by Lemma~\ref{lem:P}. Its running time is $O(\nu(\log\nu)^{s/2})=o(n^2)$, as in the proof of
Proposition~\ref{prop:case2}. So $A\in\parTIME(t'')$ with $t''=o(n^2)$, which contradicts the choice of $A$.
\end{proof}

\begin{proof}[Proof of Theorem~\ref{thm:Uplus}]
If characteristic $2$ has no ports, $L=\TRANS$ is in $\coNTIME[n;\lfloor\log n\rfloor+1]\subseteq\coNTIME[n;2\log n]$ and not in
$\NTIME(n(\log n)^{1/4})$ (Corollary~\ref{cor:A}). If it has ports, $L=A''$ of Lemma~\ref{lem:B} is in $\coNTIME[n;2\log n]$ and not in
$\parTIME(n(\log n)^{s/2})$. In both cases $L$ is not in the intersection. Finally
$\mathrm{UTIME}(f)\subseteq\NTIME(f)\cap\parTIME(f)$.
\end{proof}

\begin{remark}[Limits of the method]
(1) \emph{No separation of $\NTIME$ from $\coNTIME$.} The proof of Lemma~\ref{lem:B} does not work for NTIME machines: for a
machine $P$ that may have several accepting computations, $\sum_u\acc(P,w_u)$ modulo $2$ need not equal $\#\{u:w_u\in A''\}$
modulo $2$. So
$\coNTIME[n;2\log n]\not\subseteq\NTIME(n(\log n)^{\sigma})$ is obtained only when characteristic $2$ has no ports.
(2) \emph{Parity machines in Case~1.} Lemma~\ref{lem:A} fails for parity machines, on the reject side: on $(X,Y)$ with $Y\ne X^T$ a
parity machine may have a positive even number of accepting computations, and the part with record $R$ may be odd. For example, a
machine that guesses a bit $b$, and for $b=0$ reads the input from left to right and accepts (a head schedule that depends only on
$n$, never moving its second work head), for $b=1$ first moves its second work head $N$ cells to the right during the first epoch
(so at the first epoch boundary that head is in another block, and the two branches have different records) and then decides
$\TRANS$ and accepts iff the input is not in $\TRANS$, has one accepting computation on each input in $\TRANS$ and two on every other input; the only record class
with an odd count on $(X,X^T)$ is the $b=0$ class, whose network gives the all-ones vector in $Y$. So the record-split construction
does not, in general, give a network that satisfies Proposition~\ref{prop:table}. (This machine is slow; it shows only that this
step has no proof.)
(3) \emph{Bounded ambiguity.} Lemma~\ref{lem:A} holds for all NTMs. The proof of Lemma~\ref{lem:B} extends to machines that
accept iff they have at least one accepting computation, that have at most $k(\nu)\le(\log\nu)^{\varepsilon}$ of them on inputs of
length $\nu$, and that run in time $O(\nu(\log\nu)^{\sigma})$, whenever $\sigma+\varepsilon<s$ (for instance $\sigma=s/2$ and
$\varepsilon<s/2$). For $0\le a\le k$ one has $[a\ge1]\equiv\sum_{j=1}^{k}\binom aj\pmod2$
(the sum is $2^a-1$), the identity behind $\mathrm{FewP}\subseteq\oplus\mathrm P$. So the parity machine, for each $u$, guesses
$j$ with $1\le j\le k$ (one binary code for each value; other codes reject) and then $j$ transcripts of the given machine on $w_u$ (each in the fixed transcript encoding of
Theorem~\ref{thm:hier}, and each cut by a clock at the given machine's time bound), keeps the previous transcript on a work tape to check that they come in strictly increasing order, checks
that each is an accepting computation, and accepts iff all checks pass; each $j$-set of accepting computations gives exactly one accepting path. This costs
$O(k)$ times the running time of the given machine, which is $o(n^2)$ when $\sigma+\varepsilon<s$. For polynomial $k$ this route needs every term $j\le k$ and the question stays open.
(4) \emph{Open:} $\NTIME[n;2\log n]\not\subseteq\mathrm{UTIME}(n(\log n)^{\sigma})$ (the complement side). In Case~2 it follows,
since $\parTIME$ is closed under complement. In Case~1 the method does not give it: a machine for the complement of $\TRANS$ gives
support on all $Y\ne X^T$, not a basis vector, and Proposition~\ref{prop:table} does not apply. (No method can show that the
complement of $\TRANS$ is outside $\NTIME(n(\log n)^{\sigma})$, since it is in $\NTIME(n)$.)
\end{remark}

%======================================================================
\section{The Lean proof: what is formalized where}\label{sec:lean}

The repository of this paper (\url{https://github.com/jackryan225/ntime-dtime-polylog}, folder \src{lean/}) has a Lean~4 proof of
Theorem~\ref{thm:main} and Theorem~\ref{thm:Uplus} with no hypothesis. This section says what it proves, what it takes from
OpenAI, and which part of this text is in which Lean files. The files \src{LEAN-STATEMENT.md} and \src{LEAN-AUDIT.md} of the
repository have the details.

\paragraph{The four theorems.} The statement files in \src{lean/C1Challenge/} import only Mathlib. They define the machines, the
running time and the classes in full. The proved theorems are:
\begin{itemize}
\item \texttt{P1Final.c1\_closed}: there are $\sigma>0$ and $L\in\NTIME[n;\lfloor2\log n\rfloor]$ with $L\notin\DTIME(n(\log n)^{\sigma})$;
\item \texttt{P1Final.c1\_ntime\_closed}: the same with $L\in\NTIME(n)$ in addition, so both sentences of Theorem~\ref{thm:main}
  hold with one $\sigma$ and one language;
\item \texttt{P1Final.c1u\_closed}: Theorem~\ref{thm:Uplus};
\item \texttt{P1Final.c1u\_utime\_closed}: its consequence for $\mathrm{UTIME}$.
\end{itemize}
Each has exactly the type of its target in the statement files (compared as Lean terms), and each depends only on the axioms
\texttt{propext}, \texttt{Classical.choice} and \texttt{Quot.sound}.

\paragraph{What was checked.} Three kernels checked the proof: Lean's kernel during the build; \texttt{leanchecker}, which
replays all compiled modules; and \texttt{lean4lean}, an independent implementation of the kernel, on an export of each theorem
with its whole closure, Mathlib included. The closure has $3{,}067$ source files: $746$ of ours, $2{,}290$ of OpenAI's library
(six of them changed by our two patches) and $31$ of the library Lemmata, which OpenAI's library uses. A build of the released
files from an empty folder passed.

\paragraph{What the Lean proof takes from OpenAI's library} (commit \texttt{3344da5} of \src{github.com/openai/math}): the machine
model (copied field for field into the statement file); the port theorem and the wire expanders; the compressor with its
shrinking and refresh substeps; the packing data and the motif blocks; the network kernel; the stack-program and code calculus;
and the integer-multiplication program library. OpenAI's library states the step from ports to wire expanders for
characteristic $0$; our two patches (\src{field-uniform.patch}, \src{port-over.patch}) make it uniform in the field. All of
OpenAI's files in the build are compiled and kernel-checked together with ours. The Lean proof does not state H1--H12 one by
one. Where this text uses an input, the Lean proof uses a theorem of OpenAI's library or a lemma of our own; where
Section~\ref{sec:hyp} says that an input rests on our reading of the preprints, the Lean proof proves what it needs, or does not
need the input.

\paragraph{Differences between the Lean statements and this text.}
\begin{enumerate}
\item \emph{Machines.} The machines of this text have one-way tapes with an end marker and the guess on a second tape. The Lean
  machines have two-way tapes and take the guess as a prefix of the input. The two models give the same classes, since each
  simulates the other in $O(n+g)$ extra steps ($g$ the number of guess bits). This equivalence is argued in words; it is not in
  Lean.
\item \emph{Constants.} The clocks and the number of guess bits at $n\le1$ differ by constants. The classes are the same.
\item \emph{The case split.} Lean splits, by excluded middle, on one proposition: whether the algebraic closure of $\F_2$ has
  ports in the sense of OpenAI's library. This text splits on ``characteristic $2$ has ports'' and uses H1(b) to pass between
  fields of characteristic $2$. The Lean proof does not need H1(b). In both forms $\sigma$ is not effective.
\end{enumerate}

\paragraph{Which part of this text is in which folder.} The table is a guide to the folders of \src{lean/}. It is not a
line-by-line correspondence.
\begin{center}
\begin{tabular}{p{0.50\textwidth}p{0.43\textwidth}}
\hline
Part of this text & Lean folders \\
\hline
Statements of Theorems~\ref{thm:main} and~\ref{thm:Uplus} & \texttt{C1Challenge/} (targets), \texttt{P1Final/} (proved theorems) \\
Counters and batched lookup (Section~\ref{sec:models}) & \texttt{P1Batch/} \\
Parity time hierarchy (Theorem~\ref{thm:hier}) & \texttt{P1Hier/} \\
Case 1, transposition (Section~\ref{sec:case1}) & \texttt{P1Case1/} \\
Skeletons, the network, access chains (Sections~\ref{subsec:normal}, \ref{subsec:network}) & \texttt{P1Net/}; Bene\v{s} routing: \texttt{P1Benes/} \\
Locality lemma (Lemma~\ref{lem:locality}) & \texttt{P1P7/} \\
Witness format and its length (Section~\ref{subsec:witness}, Lemma~\ref{lem:space}) & \texttt{P1Wit/} \\
Checks C1--C4 and C14 & \texttt{P1ChkDir/} \\
Checks C5--C9 & \texttt{P1ChkGreedy/} \\
Checks C10--C13 & \texttt{P1ChkFinal/}, \texttt{P1C10/} \\
The checks as tape programs within the time bound & \texttt{P1ProgDir/}, \texttt{P1ProgEng/}, \texttt{P1Dispatch/}, \texttt{P1P5/}, \texttt{P1Sol*/}, \texttt{P1OaiReuse/} \\
Parity bookkeeping and assembly of Lemma~P & \texttt{P1Count/}, \texttt{P1Int/} \\
Wrapper (Section~\ref{sec:wrapper}); the step from the closure of $\F_2$ to a finite algebra; $\NTIME(n)$ & \texttt{P1Wrap/}, \texttt{P1NTime/}, \texttt{P1Final/} \\
Stronger form (Section~\ref{sec:utime}) & \texttt{P1C1U/}, \texttt{P1Final/} \\
\hline
\end{tabular}
\end{center}

\paragraph{Two errors of the first version, found by the Lean work.} Both are corrected in this version. Neither affected the
theorem.
\begin{enumerate}
\item[(E1)] Check C12 said that every bit of S3 ``after position'' $L$ is $0$. With positions counted from $0$ this left the bit
  at position $L$ unchecked. Section~\ref{subsec:witness} now checks every position $i\ge L$. Three independent readings of the
  section made for the Lean work found it.
\item[(E2)] The first version gave no reason for the length of the energy prefixes. The bound is now stated in
  Section~\ref{subsec:panels}. It was found when a Lean step failed: the stored energy sums were known only modulo a power of two.
\end{enumerate}

\paragraph{What is not formalized.} The equivalence of the two machine models above; the comparison with prior work
(Section~\ref{sec:prior}); the remarks of Section~\ref{sec:remarks}; and the match between the inputs H1--H12 as restated in
Section~\ref{sec:hyp} and the OpenAI preprints, which the Lean proof does not need. No professional mathematician has refereed
the Lean statements; two blind readings of each main statement by other models are recorded in \src{LEAN-STATEMENT.md}.

%======================================================================
\section{Remarks}\label{sec:remarks}

\begin{remark}[Range of the exponent]\label{rem:range}
Write $B=\lambda^{a+o(1)}$ and $T=s\log\lambda+O(1)$. The construction needs $a>s$ (the skeleton and the per-position records, of
order $t\lambda^{-a}\operatorname{polylog}\lambda$, must be $o(t\lambda^{-s})$) and $a+2s\log C_2<1$ (exhaustive evaluation of
$W_v$ in time $t^{o(1)}$, by \eqref{eq:hidden}); the big-block exponent can be raised to keep $\lambda^{\Omega(1)}$ lanes. Hence this
proof certifies every $s<1/(1+2\log C_2)$, not only the $s$ fixed above. (This range was found in the first written version of the port case and confirmed in its line-by-line check.)
\end{remark}

\begin{remark}[An obstruction to this design]
Count the variables of the literal expansion of a vertex tensor $W_v$, before any tie is applied. Each interval factor
placed at $v$ brings its boundary half-wires (the state and content strings on its in- and out-edges, $\Theta_M(B)$ bits), which are
summed inside $v$ (H7: each whole old factor goes to one new position, with its incident indices summed there); its transition codes
are not variables, since $\Phi_i$ is evaluated as a table. There are $\Theta_M(t/B)$ interval factors and at most
$O_M(t\log\lambda/(B\kappa_0^T))$ final positions, each initial factor lying whole at one final position (H8). So some final position
receives $\Omega_M(\kappa_0^T/\log\lambda)$ interval factors, that is at least $\Omega_M(B\kappa_0^T/\log\lambda)$ raw variables,
which is at least $\lambda^{s(1+\log\kappa_0)-o(1)}$ once $B\ge\lambda^{s}$. If $s(1+\log\kappa_0)>1$ an enumeration of these raw
variables, as in C13, costs $2^{\omega(\lambda)}=t^{\omega(1)}$. So this design with this evaluation cannot give every $s$; a larger
exponent needs a cheaper evaluation of the interval parity tables (for example after simplifying the kernels) or a compressor with a
better $\kappa_0$--$C_2$ trade-off. It is not an obstruction to the theorem.
\end{remark}

\begin{remark}[What the proof does not need]
No block-respecting simulation (Hopcroft--Paul--Valiant) is used: the three-block windows give the network directly
(Section~\ref{subsec:network}). The descriptor RAM of H12 is not run; only its construction is used, and the global records it
schedules (colours, offsets, layout) are replaced by locally checkable records. The proof in this text does not use the Lean proof.
Several inputs are, at the level of statements, instances of theorems of OpenAI's Lean library for every field (H2, price
part: \texttt{wire\_cost\_le}; H5: \texttt{paired\_entropy\_at\_placement}; H6: \texttt{MotifConversion.exists\_conversion},
\src{OAI/Computability/VPVNP/MotifConversion.lean} l.327; for H8, \texttt{Compression.exists\_iterated\_compression} under
\texttt{HasWireExpanders} proves the bounds of H8 for OpenAI's own construction, which is not literally H8 with the scope used
here). In this text the other inputs, or parts of inputs, rest on our reading of the preprints: H1(b), H1(d) in characteristic $2$, the class
$\X_2$ of H2, H3 in the general-class form, H4, H7, the scope of H8 over all admissible choices, and the constructions H9--H12, whose match with
OpenAI's Lean construction files this text does not check (Remark~\ref{rem:weak}); the accompanying file HYPOTHESES.md has the table. The Lean proof does not depend on these readings (Section~\ref{sec:lean}). Weak ETH is not used: Case~1 replaces it by the no-ports hypothesis itself.
\end{remark}

\begin{remark}[Weakest points of the proof in words, for a reviewer]\label{rem:weak}
These points concern the proof in this text. The Lean proof does not use our restatements H1--H12, so the theorems do not depend
on them.
\begin{enumerate}
\item H8 over $\F_2$ (EG asserts it for every field with ports; the energy estimate and the field of the parameter solution are
  the content). The proof in words is exactly as strong as H8 in Case~2. OpenAI's Lean library proves the bounds of H8 for its
  own construction, for every field $k$ under the hypothesis \texttt{HasWireExpanders k} (\texttt{Compression.exists\_iterated\_compression},
  \src{OAI/Computability/VPVNP/IteratedCompression.lean} l.115--136; arbitrary finite base, one lane displacement per base slot,
  lane-dependent widths, energy bound $\sqrt{\mathcal L_T}\le(\sqrt{\mathcal L_0}+O(P_0))/2^T$). The step from H1 to
  \texttt{HasWireExpanders} is formalized there only for characteristic $0$ (\texttt{wireExpanders\_of\_port},
  \src{OAI/Computability/MultitapeWeakETH/PortTemplates.lean} l.45). For the first version we read but did not build these files; the Lean proof builds them, with two patches of ours that make
  the port statement uniform in the field (Section~\ref{sec:lean}).
  A separate check shows that ports in any field of characteristic $2$ give ports over $\F_2$ with the
  same graphs, and that the parameter solution lies in $\F_2$.
\item The characteristic-$2$ form H4 of T's geometric lemma, and H1(d) in characteristic $2$. They rest on a reading of the proofs
  (EG \src{port-graphs.tex} l.38--54); for the first version no Lean build existed; the field-uniform patch is now part of the Lean proof and compiles
  (Section~\ref{sec:lean}). H5 is an instance of a field-general Lean statement (H5 above).
\item Lemma~\ref{lem:locality} rests on the constructions H9--H11 as we restate them from EG's text. A reviewer should compare
  Step~1 (placement) and Step~2 (the four parent bases of a motif) with \src{EG/compressor-appendix.tex} l.39--55 and
  \src{EG/compressor.tex} l.392--404. OpenAI's Lean construction, which the Lean form of H8 certifies, fixes some free choices of
  the text in its own way: we follow its greedy scans exactly ((R1), (R2)); its slot identifiers use absolute base addresses where
  H12 writes lanes relative to a source lane, and its layout scores use the weights $(10,20,6,5)$ of the four surviving supply types.
  A proof against the Lean construction would use those keys and weights in C7--C10;
  this changes no radius of Lemma~\ref{lem:locality}. An independent reading of OpenAI's executed Lean construction (commit
  \texttt{3344da5}), datum by datum, found that H9--H11 hold as used here, with the same radii and the same uniqueness argument; the
  only differences are these keys and weights (\src{CHECKS.md}, row 28).
\item Lemma~\ref{lem:pathcount}: the access chains of Section~\ref{subsec:network}(c) and the parity of the Bene\v{s} routing
  (settings are computed, not guessed; a guessed setting would multiply each term by the number of valid settings, an even number
  for most permutations, and destroy the parity).
\item The uniqueness of the passing witness (Lemma~\ref{lem:unique}): every free choice of EG must be fixed by a rule that the
  checks enforce; we listed them under ``Fixed choices''.
\end{enumerate}
\end{remark}

\begin{remark}[Why characteristic $2$]
We use characteristic $2$ because then every vertex tensor $W_v$ of \eqref{eq:expansion} is $0$ or $1$, so a witness needs no
multiplicity bits. The argument extends to every prime $p$. For $\mathrm{Mod}_p\mathrm{TIME}$ (accept iff
$\acc\not\equiv0\pmod p$) there is a time hierarchy as well: the class is closed under complement with a constant factor in time,
since a machine with $1+(p-1)\acc(M,x)^{p-1}$ accepting computations ($p-1$ sequential copies of $M$, then Fermat's little
theorem) accepts exactly when $\acc(M,x)\equiv0$, and the diagonal machine of Theorem~\ref{thm:hier} uses this complement in place
of the extra accepting computation. In Case~2 over $\F_p$ the vertex tensors lie in $\F_p$, so the witness carries, for each final
position $v$, an index $c_v<W_v$ (with $W_v$ read as an integer in $[0,p)$): $O(\log p)$ more bits per final position, which fit
into the bound of Lemma~\ref{lem:space}. This needs the characteristic-$p$ forms of H1(d), H2, H4, H5 and H8, and gives
Theorem~\ref{thm:main} again; Section~\ref{sec:utime} would need further changes. We do not pursue it.
\end{remark}

\begin{remark}[The number of guess bits]
The proof gives slightly more than stated: Case~1 uses $\lfloor\log n\rfloor+1$ guess bits and Case~2 uses
$\lceil\log t\rceil\le\log\nu+\log\log\nu+2$ guess bits at input length $\nu$ (for $n\ge n_M$, since
$\nu\ge m\ge t\lambda^{-s}$). So $\NTIME[n;\,g]\not\subseteq\DTIME(n(\log n)^{\sigma})$ already for
$g(n)=\lceil\log n\rceil+\lceil\log\log n\rceil+2$, which is below $\lfloor2\log n\rfloor$ for all large $n$; $\lfloor2\log n\rfloor$ is a
safe, weaker choice.
\end{remark}

\begin{remark}[What is not claimed]
(1) No explicit $\sigma$: $\sigma=1/4$ if characteristic $2$ has no ports and $\sigma=s/2$ if it has ports, with $s$ from constants
that H8 asserts to exist. (2) No separation of $\NTIME$ from
$\coNTIME$ (Section~\ref{sec:utime}). (3) Nothing for random-access machines: there $\TRANS$ is decidable in linear time, so Case~1
has no analogue. (4) The argument does not relativize: it uses that one step of a multitape machine reads only the cells under its
heads, so that a short time interval touches $O(1)$ small blocks. For oracle machines with a query tape that is erased after each
query (queries of linear length), relative to a set $A$ that is complete for linear space under linear-time reductions,
$\NTIME^{A}(n)=\DTIME^{A}(O(n))$ (simulate the nondeterministic oracle machine in linear space, then make one query; Hemmerling~\cite[Proposition 15]{Hem}, \cite[Corollary 1]{Hem13}). So any proof of Theorem~\ref{thm:main} must fail to relativize.
\end{remark}

\begin{thebibliography}{99}
\bibitem{EG} OpenAI, \emph{Euclidean geometry in the asymptotic spectrum of tensors}, preprint, 24 November 2026,
  \texttt{github.com/openai/math}, folder \url{preprints/Euclidean-geometry-in-the-asymptotic-spectrum-of-tensors-November-24-2026}.
\bibitem{CL} OpenAI, \emph{Circuit lower bounds for Hamiltonian cycle and the permanent}, preprint, 24 November 2026,
  \texttt{github.com/openai/math}, folder \url{preprints/Circuit-lower-bounds-for-Hamiltonian-cycle-and-the-permanent-November-24-2026}.
\bibitem{T} OpenAI, \emph{The logarithmic exponent of integer multiplication on tapes}, preprint, 24 November 2026,
  \texttt{github.com/openai/math}, folder
  \url{preprints/The-logarithmic-exponent-of-integer-multiplication-on-tapes-November-24-2026}. Lean library: same repository,
  folder \texttt{lean} (commit \texttt{3344da5}).
\bibitem{OAI-ETH} OpenAI, \emph{The Exponential Time Hypothesis is False}, preprint, 25 October 2026, \texttt{github.com/openai/math}.
\bibitem{Ant} Claude (Anthropic), \emph{VP $\ne$ VNP over every field}, hexagonmath.org 2611.00954v1, 23 November 2026.
\bibitem{PPST} W.~J.~Paul, N.~Pippenger, E.~Szemer\'edi, W.~T.~Trotter, \emph{On determinism versus non-determinism and related
  problems}, 24th FOCS (1983), 429--438.
\bibitem{San} R.~Santhanam, \emph{On separators, segregators and time versus space}, 16th CCC (2001), 286--294.
\bibitem{CW} L.~Chen, Y.~Wang, \emph{Separating RAM and multitape Turing machines with short random oracles}, ECCC TR26-006 (2026).
\bibitem{Wil} R.~Williams, \emph{Inductive time-space lower bounds for SAT and related problems}, Computational Complexity 15 (2006),
  433--470.
\bibitem{Lin} A.~Lingas, \emph{Consequences of APSP, triangle detection, and 3SUM hardness for separation between determinism and
  non-determinism}, arXiv:2011.02428v5 (2021); Procedia Computer Science 195 (2021), 163--171.
\bibitem{AVW} J.~Alman, V.~Vassilevska Williams, \emph{Truly subquadratic 3SUM and truly subcubic APSP via triangles in sparse
  lopsided graphs}, arXiv:2610.06783 (2026).
\bibitem{HvdH} D.~Harvey, J.~van der Hoeven, \emph{Integer multiplication is at least as hard as matrix transposition}, FOCS 2025,
  1366--1381; arXiv:2503.22848.
\bibitem{Tal} Y.~Tal, \emph{Sorting in $o(N\log N)$ time on a multitape Turing machine}, ECCC TR26-331 (2026).
\bibitem{SS} C.~Strassle, B.~Subercaseaux, \emph{Simulating nondeterministic Turing machines in subexponential deterministic time},
arXiv:2611.11832 (9 November 2026).
\bibitem{BGW} R.~V.~Book, S.~A.~Greibach, B.~Wegbreit, \emph{Time- and tape-bounded Turing acceptors and AFLs}, J. Comput. System
  Sci. 4 (1970), 606--621.
\bibitem{HS65} J.~Hartmanis, R.~E.~Stearns, \emph{On the computational complexity of algorithms}, Trans. Amer. Math. Soc. 117
  (1965), 285--306.
\bibitem{Cook73} S.~A.~Cook, \emph{A hierarchy for nondeterministic time complexity}, J. Comput. System Sci. 7 (1973), 343--353.
\bibitem{SFM78} J.~I.~Seiferas, M.~J.~Fischer, A.~R.~Meyer, \emph{Separating nondeterministic time complexity classes}, J. ACM 25
  (1978), 146--167.
\bibitem{Zak83} S.~\v{Z}\'ak, \emph{A Turing machine time hierarchy}, Theoret. Comput. Sci. 26 (1983), 327--333.
\bibitem{Benes} V.~E.~Bene\v{s}, \emph{Optimal rearrangeable multistage connecting networks}, Bell System Tech. J. 43 (1964),
  1641--1656.
\bibitem{Waksman} A.~Waksman, \emph{A permutation network}, J. ACM 15 (1968), 159--163.
\bibitem{HPV} J.~Hopcroft, W.~Paul, L.~Valiant, \emph{On time versus space}, J. ACM 24 (1977), 332--337.
\bibitem{BG} R.~Beigel, J.~Gill, \emph{Counting classes: thresholds, parity, mods, and fewness}, Theoret. Comput. Sci. 103 (1992),
  3--23.
\bibitem{Hem} A.~Hemmerling, \emph{Observations on complete sets between linear time and polynomial time}, Inform. and Comput. 209
  (2011), 173--182.
\bibitem{Hem13} A.~Hemmerling, \emph{On the tape-number problem for deterministic time classes}, arXiv:1309.7457 (2013).
\end{thebibliography}

\end{document}
